\documentclass[11pt,a4paper]{article}
% =====================================================================
%  Preámbulo común para los apuntes de Física Electrónica
%  Uso (desde unidades/XX-nombre/archivo.tex):
%     \documentclass[11pt,a4paper]{article}
%     % =====================================================================
%  Preámbulo común para los apuntes de Física Electrónica
%  Uso (desde unidades/XX-nombre/archivo.tex):
%     \documentclass[11pt,a4paper]{article}
%     % =====================================================================
%  Preámbulo común para los apuntes de Física Electrónica
%  Uso (desde unidades/XX-nombre/archivo.tex):
%     \documentclass[11pt,a4paper]{article}
%     \input{../../comun/preambulo.tex}
% =====================================================================

% ---------- Idioma y tipografía ----------
\usepackage[T1]{fontenc}
\usepackage[utf8]{inputenc}
\usepackage[spanish,es-nodecimaldot,es-noshorthands,es-tabla]{babel}
\usepackage{lmodern}
\usepackage{microtype}
\usepackage[a4paper,margin=2.3cm,headheight=14pt]{geometry}
\usepackage{parskip}
\setlength{\parskip}{0.55em}

% ---------- Matemática y unidades ----------
\usepackage{amsmath,amssymb,mathtools}
\usepackage{bm}
\usepackage[per-mode=symbol,output-decimal-marker={.},separate-uncertainty]{siunitx}
\sisetup{range-phrase={\text{--}},range-units=single}
\DeclareSIUnit{\eV}{eV}

% ---------- Gráficos ----------
\usepackage{graphicx}
\usepackage{xcolor}
\usepackage{tikz}
\usetikzlibrary{arrows.meta,decorations.pathmorphing,decorations.markings,
  patterns,calc,positioning,shapes.geometric,fit,backgrounds,shadings}
\usepackage{pgfplots}
\pgfplotsset{compat=1.18}
\usepgfplotslibrary{fillbetween}
\usepackage[siunitx,RPvoltages]{circuitikz}

% ---------- Tablas, listas, figuras ----------
\usepackage{booktabs,array,tabularx,multirow}
\usepackage{enumitem}
\setlist{itemsep=0.2em,topsep=0.3em}
\usepackage{caption,subcaption}
\captionsetup{font=small,labelfont=bf,format=hang}
\usepackage{float}
\usepackage{wrapfig}
\usepackage{needspace}

% ---------- Colores ----------
\definecolor{azul}{HTML}{1F4E79}
\definecolor{azulclaro}{HTML}{DCE9F5}
\definecolor{verde}{HTML}{2E7D32}
\definecolor{verdeclaro}{HTML}{E3F2E1}
\definecolor{naranja}{HTML}{C55A11}
\definecolor{naranjaclaro}{HTML}{FBE9DB}
\definecolor{rojo}{HTML}{B71C1C}
\definecolor{rojoclaro}{HTML}{FDECEA}
\definecolor{violeta}{HTML}{5E35B1}
\definecolor{violetaclaro}{HTML}{EDE7F6}
\definecolor{gris}{HTML}{555555}
% colores para dibujos de bandas / portadores
\colorlet{colEc}{azul}
\colorlet{colEv}{rojo}
\colorlet{colElec}{azul}
\colorlet{colHueco}{rojo}
\colorlet{colFoton}{naranja}

% ---------- Cajas didácticas ----------
\usepackage[most]{tcolorbox}
\tcbset{
  base/.style={enhanced,breakable,arc=2pt,boxrule=0.6pt,left=6pt,right=6pt,
    top=4pt,bottom=4pt,fonttitle=\bfseries\small,before skip=8pt,after skip=8pt}
}
% Definición
\newtcolorbox{definicion}[1][]{base,colback=azulclaro,colframe=azul,
  title={Definición\ifx\relax#1\relax\else: #1\fi}}
% Resultado / fórmula clave
\newtcolorbox{clave}[1][]{base,colback=verdeclaro,colframe=verde,
  title={\ifx\relax#1\relax Resultado clave\else #1\fi}}
% Ejemplo resuelto
\newtcolorbox{ejemplo}[1][]{base,colback=white,colframe=violeta,
  coltitle=white,colbacktitle=violeta,
  title={Ejemplo\ifx\relax#1\relax\else: #1\fi}}
% Corrección / aclaración respecto de las diapositivas
\newtcolorbox{ojo}[1][]{base,colback=rojoclaro,colframe=rojo,
  title={\ensuremath{\triangle}\ Ojo\ifx\relax#1\relax\else: #1\fi}}
% Intuición física / comentario
\newtcolorbox{intuicion}[1][]{base,colback=naranjaclaro,colframe=naranja,
  title={Idea física\ifx\relax#1\relax\else: #1\fi}}

% ---------- Encabezados ----------
\usepackage{fancyhdr}
\pagestyle{fancy}
\fancyhf{}
\renewcommand{\headrulewidth}{0.4pt}
\fancyhead[L]{\small\color{gris}Física Electrónica}
\fancyhead[R]{\small\color{gris}\unidadtitulo}
\fancyfoot[C]{\small\thepage}
\newcommand{\unidadtitulo}{}
\newcommand{\unidad}[1]{\renewcommand{\unidadtitulo}{#1}}

\usepackage{titlesec}
\titleformat{\section}{\Large\bfseries\color{azul}}{\thesection}{0.6em}{}
\titleformat{\subsection}{\large\bfseries\color{azul}}{\thesubsection}{0.6em}{}
\titleformat{\subsubsection}{\normalsize\bfseries\color{azul!80!black}}{\thesubsubsection}{0.6em}{}

\usepackage[hidelinks,bookmarksnumbered]{hyperref}
\usepackage{bookmark}

% ---------- Macros físicas ----------
\newcommand{\EG}{E_G}
\newcommand{\Ec}{E_c}
\newcommand{\Ev}{E_v}
\newcommand{\EF}{E_F}
\newcommand{\Eph}{E_{\mathrm{ph}}}
\newcommand{\lG}{\lambda_G}
\newcommand{\VT}{V_T}
\newcommand{\Popt}{P_{\mathrm{opt}}}
\newcommand{\Resp}{\mathcal{R}}
\newcommand{\vsat}{v_{\mathrm{sat}}}
\newcommand{\Voc}{V_{oc}}
\newcommand{\Isc}{I_{sc}}
\newcommand{\diapo}[1]{{\small\textsf{[diapo~#1]}}}
\newcommand{\dd}{\,\mathrm{d}}

% Distintivo para resultados que ya están en la hoja de fórmulas oficial.
% Uso: en el título de una caja, \begin{clave}[Título \enhoja], o en texto con \enhojatexto.
\newcommand{\enhojatexto}{{\scriptsize\sffamily\colorbox{verdeclaro}{\textcolor{verde!70!black}{\ensuremath{\checkmark}\,en la hoja}}}}
\newcommand{\enhoja}{\hfill{\scriptsize\sffamily\colorbox{white}{\textcolor{verde!70!black}{\ensuremath{\checkmark}\,en la hoja}}}}

\input{../../comun/constantes.tex}

% ---------- Estilos TikZ comunes ----------
\tikzset{
  banda/.style={line width=1.1pt},
  bandac/.style={banda,colEc},
  bandav/.style={banda,colEv},
  fermi/.style={dashed,gris,line width=0.6pt},
  elec/.style={circle,fill=colElec,inner sep=0pt,minimum size=5pt},
  hueco/.style={circle,draw=colHueco,fill=white,line width=0.9pt,inner sep=0pt,minimum size=5pt},
  foton/.style={colFoton,line width=0.9pt,decorate,
    decoration={snake,amplitude=1.6pt,segment length=5pt,post length=3pt},-{Stealth[length=5pt]}},
  flecha/.style={-{Stealth[length=5pt]},line width=0.7pt},
}

% =====================================================================

% ---------- Idioma y tipografía ----------
\usepackage[T1]{fontenc}
\usepackage[utf8]{inputenc}
\usepackage[spanish,es-nodecimaldot,es-noshorthands,es-tabla]{babel}
\usepackage{lmodern}
\usepackage{microtype}
\usepackage[a4paper,margin=2.3cm,headheight=14pt]{geometry}
\usepackage{parskip}
\setlength{\parskip}{0.55em}

% ---------- Matemática y unidades ----------
\usepackage{amsmath,amssymb,mathtools}
\usepackage{bm}
\usepackage[per-mode=symbol,output-decimal-marker={.},separate-uncertainty]{siunitx}
\sisetup{range-phrase={\text{--}},range-units=single}
\DeclareSIUnit{\eV}{eV}

% ---------- Gráficos ----------
\usepackage{graphicx}
\usepackage{xcolor}
\usepackage{tikz}
\usetikzlibrary{arrows.meta,decorations.pathmorphing,decorations.markings,
  patterns,calc,positioning,shapes.geometric,fit,backgrounds,shadings}
\usepackage{pgfplots}
\pgfplotsset{compat=1.18}
\usepgfplotslibrary{fillbetween}
\usepackage[siunitx,RPvoltages]{circuitikz}

% ---------- Tablas, listas, figuras ----------
\usepackage{booktabs,array,tabularx,multirow}
\usepackage{enumitem}
\setlist{itemsep=0.2em,topsep=0.3em}
\usepackage{caption,subcaption}
\captionsetup{font=small,labelfont=bf,format=hang}
\usepackage{float}
\usepackage{wrapfig}
\usepackage{needspace}

% ---------- Colores ----------
\definecolor{azul}{HTML}{1F4E79}
\definecolor{azulclaro}{HTML}{DCE9F5}
\definecolor{verde}{HTML}{2E7D32}
\definecolor{verdeclaro}{HTML}{E3F2E1}
\definecolor{naranja}{HTML}{C55A11}
\definecolor{naranjaclaro}{HTML}{FBE9DB}
\definecolor{rojo}{HTML}{B71C1C}
\definecolor{rojoclaro}{HTML}{FDECEA}
\definecolor{violeta}{HTML}{5E35B1}
\definecolor{violetaclaro}{HTML}{EDE7F6}
\definecolor{gris}{HTML}{555555}
% colores para dibujos de bandas / portadores
\colorlet{colEc}{azul}
\colorlet{colEv}{rojo}
\colorlet{colElec}{azul}
\colorlet{colHueco}{rojo}
\colorlet{colFoton}{naranja}

% ---------- Cajas didácticas ----------
\usepackage[most]{tcolorbox}
\tcbset{
  base/.style={enhanced,breakable,arc=2pt,boxrule=0.6pt,left=6pt,right=6pt,
    top=4pt,bottom=4pt,fonttitle=\bfseries\small,before skip=8pt,after skip=8pt}
}
% Definición
\newtcolorbox{definicion}[1][]{base,colback=azulclaro,colframe=azul,
  title={Definición\ifx\relax#1\relax\else: #1\fi}}
% Resultado / fórmula clave
\newtcolorbox{clave}[1][]{base,colback=verdeclaro,colframe=verde,
  title={\ifx\relax#1\relax Resultado clave\else #1\fi}}
% Ejemplo resuelto
\newtcolorbox{ejemplo}[1][]{base,colback=white,colframe=violeta,
  coltitle=white,colbacktitle=violeta,
  title={Ejemplo\ifx\relax#1\relax\else: #1\fi}}
% Corrección / aclaración respecto de las diapositivas
\newtcolorbox{ojo}[1][]{base,colback=rojoclaro,colframe=rojo,
  title={\ensuremath{\triangle}\ Ojo\ifx\relax#1\relax\else: #1\fi}}
% Intuición física / comentario
\newtcolorbox{intuicion}[1][]{base,colback=naranjaclaro,colframe=naranja,
  title={Idea física\ifx\relax#1\relax\else: #1\fi}}

% ---------- Encabezados ----------
\usepackage{fancyhdr}
\pagestyle{fancy}
\fancyhf{}
\renewcommand{\headrulewidth}{0.4pt}
\fancyhead[L]{\small\color{gris}Física Electrónica}
\fancyhead[R]{\small\color{gris}\unidadtitulo}
\fancyfoot[C]{\small\thepage}
\newcommand{\unidadtitulo}{}
\newcommand{\unidad}[1]{\renewcommand{\unidadtitulo}{#1}}

\usepackage{titlesec}
\titleformat{\section}{\Large\bfseries\color{azul}}{\thesection}{0.6em}{}
\titleformat{\subsection}{\large\bfseries\color{azul}}{\thesubsection}{0.6em}{}
\titleformat{\subsubsection}{\normalsize\bfseries\color{azul!80!black}}{\thesubsubsection}{0.6em}{}

\usepackage[hidelinks,bookmarksnumbered]{hyperref}
\usepackage{bookmark}

% ---------- Macros físicas ----------
\newcommand{\EG}{E_G}
\newcommand{\Ec}{E_c}
\newcommand{\Ev}{E_v}
\newcommand{\EF}{E_F}
\newcommand{\Eph}{E_{\mathrm{ph}}}
\newcommand{\lG}{\lambda_G}
\newcommand{\VT}{V_T}
\newcommand{\Popt}{P_{\mathrm{opt}}}
\newcommand{\Resp}{\mathcal{R}}
\newcommand{\vsat}{v_{\mathrm{sat}}}
\newcommand{\Voc}{V_{oc}}
\newcommand{\Isc}{I_{sc}}
\newcommand{\diapo}[1]{{\small\textsf{[diapo~#1]}}}
\newcommand{\dd}{\,\mathrm{d}}

% Distintivo para resultados que ya están en la hoja de fórmulas oficial.
% Uso: en el título de una caja, \begin{clave}[Título \enhoja], o en texto con \enhojatexto.
\newcommand{\enhojatexto}{{\scriptsize\sffamily\colorbox{verdeclaro}{\textcolor{verde!70!black}{\ensuremath{\checkmark}\,en la hoja}}}}
\newcommand{\enhoja}{\hfill{\scriptsize\sffamily\colorbox{white}{\textcolor{verde!70!black}{\ensuremath{\checkmark}\,en la hoja}}}}

% =====================================================================
%  Constantes oficiales de la materia ("Tabla de constantes.pdf").
%  Usar SIEMPRE estos valores en los ejemplos numéricos, para que
%  coincidan con los que se usan en clase y en el examen.
% =====================================================================

% ---------- Constantes físicas ----------
\newcommand{\qval}{\num{1.60e-19}}          % C
\newcommand{\kval}{\num{8.61e-5}}           % eV/K
\newcommand{\kJval}{\num{1.38e-23}}         % J/K
\newcommand{\hval}{\num{6.62e-34}}          % J s
\newcommand{\heVval}{\num{4.13e-15}}        % eV s
\newcommand{\epsval}{\num{8.85e-14}}        % F/cm
\newcommand{\VTval}{\SI{26}{mV}}            % kT/q a 300 K (0.0259 V)
\newcommand{\epsrox}{3.9}                   % SiO2

% ---------- Silicio ----------
\newcommand{\epsrSi}{11.8}
\newcommand{\niSi}{\num{1.0e10}}            % cm^-3
\newcommand{\munSi}{1350}                   % cm^2/(V s)
\newcommand{\mupSi}{480}                    % cm^2/(V s)
\newcommand{\EGSi}{1.12}                    % eV
\newcommand{\NCSi}{\num{3.22e19}}           % cm^-3
\newcommand{\NVSi}{\num{1.83e19}}           % cm^-3

% ---------- Germanio ----------
\newcommand{\epsrGe}{15.8}
\newcommand{\niGe}{\num{2.4e13}}
\newcommand{\munGe}{3900}
\newcommand{\mupGe}{1900}
\newcommand{\EGGe}{0.66}

% ---------- GaAs ----------
\newcommand{\epsrGaAs}{13.1}
\newcommand{\niGaAs}{\num{2.1e6}}
\newcommand{\munGaAs}{8800}
\newcommand{\mupGaAs}{400}
\newcommand{\EGGaAs}{1.42}

% ---------- Tabla para incluir en los apuntes ----------
\newcommand{\tablaconstantes}{%
\begin{tabular}{lccc}
\toprule
 & Si & Ge & GaAs \\
\midrule
$\epsilon_r$                     & 11.8 & 15.8 & 13.1 \\
$\mu_n$ (\si{cm^2/V.s})          & 1350 & 3900 & 8800 \\
$\mu_p$ (\si{cm^2/V.s})          & 480  & 1900 & 400 \\
$n_i$ (\si{cm^{-3}}, 300\,K)     & \num{1.0e10} & \num{2.4e13} & \num{2.1e6} \\
$E_G$ (\si{\eV})                 & 1.12 & 0.66 & 1.42 \\
\bottomrule
\end{tabular}}


% ---------- Estilos TikZ comunes ----------
\tikzset{
  banda/.style={line width=1.1pt},
  bandac/.style={banda,colEc},
  bandav/.style={banda,colEv},
  fermi/.style={dashed,gris,line width=0.6pt},
  elec/.style={circle,fill=colElec,inner sep=0pt,minimum size=5pt},
  hueco/.style={circle,draw=colHueco,fill=white,line width=0.9pt,inner sep=0pt,minimum size=5pt},
  foton/.style={colFoton,line width=0.9pt,decorate,
    decoration={snake,amplitude=1.6pt,segment length=5pt,post length=3pt},-{Stealth[length=5pt]}},
  flecha/.style={-{Stealth[length=5pt]},line width=0.7pt},
}

% =====================================================================

% ---------- Idioma y tipografía ----------
\usepackage[T1]{fontenc}
\usepackage[utf8]{inputenc}
\usepackage[spanish,es-nodecimaldot,es-noshorthands,es-tabla]{babel}
\usepackage{lmodern}
\usepackage{microtype}
\usepackage[a4paper,margin=2.3cm,headheight=14pt]{geometry}
\usepackage{parskip}
\setlength{\parskip}{0.55em}

% ---------- Matemática y unidades ----------
\usepackage{amsmath,amssymb,mathtools}
\usepackage{bm}
\usepackage[per-mode=symbol,output-decimal-marker={.},separate-uncertainty]{siunitx}
\sisetup{range-phrase={\text{--}},range-units=single}
\DeclareSIUnit{\eV}{eV}

% ---------- Gráficos ----------
\usepackage{graphicx}
\usepackage{xcolor}
\usepackage{tikz}
\usetikzlibrary{arrows.meta,decorations.pathmorphing,decorations.markings,
  patterns,calc,positioning,shapes.geometric,fit,backgrounds,shadings}
\usepackage{pgfplots}
\pgfplotsset{compat=1.18}
\usepgfplotslibrary{fillbetween}
\usepackage[siunitx,RPvoltages]{circuitikz}

% ---------- Tablas, listas, figuras ----------
\usepackage{booktabs,array,tabularx,multirow}
\usepackage{enumitem}
\setlist{itemsep=0.2em,topsep=0.3em}
\usepackage{caption,subcaption}
\captionsetup{font=small,labelfont=bf,format=hang}
\usepackage{float}
\usepackage{wrapfig}
\usepackage{needspace}

% ---------- Colores ----------
\definecolor{azul}{HTML}{1F4E79}
\definecolor{azulclaro}{HTML}{DCE9F5}
\definecolor{verde}{HTML}{2E7D32}
\definecolor{verdeclaro}{HTML}{E3F2E1}
\definecolor{naranja}{HTML}{C55A11}
\definecolor{naranjaclaro}{HTML}{FBE9DB}
\definecolor{rojo}{HTML}{B71C1C}
\definecolor{rojoclaro}{HTML}{FDECEA}
\definecolor{violeta}{HTML}{5E35B1}
\definecolor{violetaclaro}{HTML}{EDE7F6}
\definecolor{gris}{HTML}{555555}
% colores para dibujos de bandas / portadores
\colorlet{colEc}{azul}
\colorlet{colEv}{rojo}
\colorlet{colElec}{azul}
\colorlet{colHueco}{rojo}
\colorlet{colFoton}{naranja}

% ---------- Cajas didácticas ----------
\usepackage[most]{tcolorbox}
\tcbset{
  base/.style={enhanced,breakable,arc=2pt,boxrule=0.6pt,left=6pt,right=6pt,
    top=4pt,bottom=4pt,fonttitle=\bfseries\small,before skip=8pt,after skip=8pt}
}
% Definición
\newtcolorbox{definicion}[1][]{base,colback=azulclaro,colframe=azul,
  title={Definición\ifx\relax#1\relax\else: #1\fi}}
% Resultado / fórmula clave
\newtcolorbox{clave}[1][]{base,colback=verdeclaro,colframe=verde,
  title={\ifx\relax#1\relax Resultado clave\else #1\fi}}
% Ejemplo resuelto
\newtcolorbox{ejemplo}[1][]{base,colback=white,colframe=violeta,
  coltitle=white,colbacktitle=violeta,
  title={Ejemplo\ifx\relax#1\relax\else: #1\fi}}
% Corrección / aclaración respecto de las diapositivas
\newtcolorbox{ojo}[1][]{base,colback=rojoclaro,colframe=rojo,
  title={\ensuremath{\triangle}\ Ojo\ifx\relax#1\relax\else: #1\fi}}
% Intuición física / comentario
\newtcolorbox{intuicion}[1][]{base,colback=naranjaclaro,colframe=naranja,
  title={Idea física\ifx\relax#1\relax\else: #1\fi}}

% ---------- Encabezados ----------
\usepackage{fancyhdr}
\pagestyle{fancy}
\fancyhf{}
\renewcommand{\headrulewidth}{0.4pt}
\fancyhead[L]{\small\color{gris}Física Electrónica}
\fancyhead[R]{\small\color{gris}\unidadtitulo}
\fancyfoot[C]{\small\thepage}
\newcommand{\unidadtitulo}{}
\newcommand{\unidad}[1]{\renewcommand{\unidadtitulo}{#1}}

\usepackage{titlesec}
\titleformat{\section}{\Large\bfseries\color{azul}}{\thesection}{0.6em}{}
\titleformat{\subsection}{\large\bfseries\color{azul}}{\thesubsection}{0.6em}{}
\titleformat{\subsubsection}{\normalsize\bfseries\color{azul!80!black}}{\thesubsubsection}{0.6em}{}

\usepackage[hidelinks,bookmarksnumbered]{hyperref}
\usepackage{bookmark}

% ---------- Macros físicas ----------
\newcommand{\EG}{E_G}
\newcommand{\Ec}{E_c}
\newcommand{\Ev}{E_v}
\newcommand{\EF}{E_F}
\newcommand{\Eph}{E_{\mathrm{ph}}}
\newcommand{\lG}{\lambda_G}
\newcommand{\VT}{V_T}
\newcommand{\Popt}{P_{\mathrm{opt}}}
\newcommand{\Resp}{\mathcal{R}}
\newcommand{\vsat}{v_{\mathrm{sat}}}
\newcommand{\Voc}{V_{oc}}
\newcommand{\Isc}{I_{sc}}
\newcommand{\diapo}[1]{{\small\textsf{[diapo~#1]}}}
\newcommand{\dd}{\,\mathrm{d}}

% Distintivo para resultados que ya están en la hoja de fórmulas oficial.
% Uso: en el título de una caja, \begin{clave}[Título \enhoja], o en texto con \enhojatexto.
\newcommand{\enhojatexto}{{\scriptsize\sffamily\colorbox{verdeclaro}{\textcolor{verde!70!black}{\ensuremath{\checkmark}\,en la hoja}}}}
\newcommand{\enhoja}{\hfill{\scriptsize\sffamily\colorbox{white}{\textcolor{verde!70!black}{\ensuremath{\checkmark}\,en la hoja}}}}

% =====================================================================
%  Constantes oficiales de la materia ("Tabla de constantes.pdf").
%  Usar SIEMPRE estos valores en los ejemplos numéricos, para que
%  coincidan con los que se usan en clase y en el examen.
% =====================================================================

% ---------- Constantes físicas ----------
\newcommand{\qval}{\num{1.60e-19}}          % C
\newcommand{\kval}{\num{8.61e-5}}           % eV/K
\newcommand{\kJval}{\num{1.38e-23}}         % J/K
\newcommand{\hval}{\num{6.62e-34}}          % J s
\newcommand{\heVval}{\num{4.13e-15}}        % eV s
\newcommand{\epsval}{\num{8.85e-14}}        % F/cm
\newcommand{\VTval}{\SI{26}{mV}}            % kT/q a 300 K (0.0259 V)
\newcommand{\epsrox}{3.9}                   % SiO2

% ---------- Silicio ----------
\newcommand{\epsrSi}{11.8}
\newcommand{\niSi}{\num{1.0e10}}            % cm^-3
\newcommand{\munSi}{1350}                   % cm^2/(V s)
\newcommand{\mupSi}{480}                    % cm^2/(V s)
\newcommand{\EGSi}{1.12}                    % eV
\newcommand{\NCSi}{\num{3.22e19}}           % cm^-3
\newcommand{\NVSi}{\num{1.83e19}}           % cm^-3

% ---------- Germanio ----------
\newcommand{\epsrGe}{15.8}
\newcommand{\niGe}{\num{2.4e13}}
\newcommand{\munGe}{3900}
\newcommand{\mupGe}{1900}
\newcommand{\EGGe}{0.66}

% ---------- GaAs ----------
\newcommand{\epsrGaAs}{13.1}
\newcommand{\niGaAs}{\num{2.1e6}}
\newcommand{\munGaAs}{8800}
\newcommand{\mupGaAs}{400}
\newcommand{\EGGaAs}{1.42}

% ---------- Tabla para incluir en los apuntes ----------
\newcommand{\tablaconstantes}{%
\begin{tabular}{lccc}
\toprule
 & Si & Ge & GaAs \\
\midrule
$\epsilon_r$                     & 11.8 & 15.8 & 13.1 \\
$\mu_n$ (\si{cm^2/V.s})          & 1350 & 3900 & 8800 \\
$\mu_p$ (\si{cm^2/V.s})          & 480  & 1900 & 400 \\
$n_i$ (\si{cm^{-3}}, 300\,K)     & \num{1.0e10} & \num{2.4e13} & \num{2.1e6} \\
$E_G$ (\si{\eV})                 & 1.12 & 0.66 & 1.42 \\
\bottomrule
\end{tabular}}


% ---------- Estilos TikZ comunes ----------
\tikzset{
  banda/.style={line width=1.1pt},
  bandac/.style={banda,colEc},
  bandav/.style={banda,colEv},
  fermi/.style={dashed,gris,line width=0.6pt},
  elec/.style={circle,fill=colElec,inner sep=0pt,minimum size=5pt},
  hueco/.style={circle,draw=colHueco,fill=white,line width=0.9pt,inner sep=0pt,minimum size=5pt},
  foton/.style={colFoton,line width=0.9pt,decorate,
    decoration={snake,amplitude=1.6pt,segment length=5pt,post length=3pt},-{Stealth[length=5pt]}},
  flecha/.style={-{Stealth[length=5pt]},line width=0.7pt},
}

\unidad{Unidad 6 -- Optoelectrónica}

\title{\color{azul}\textbf{Optoelectrónica}\\[0.3em]
  \large Fotogeneración, fotoconductores, fotodiodos PN y PIN, celdas solares y LEDs}
\author{Física Electrónica -- Resumen de la unidad 6}
\date{Basado en las clases (Ing.~E.~D.~Baez) y en Pierret, \emph{Semiconductor Device Fundamentals}, cap.~3.3 y cap.~9}

\begin{document}
\maketitle
\thispagestyle{empty}

\begin{tcolorbox}[base,colback=white,colframe=gris,title={Cómo leer este apunte}]
\small
\begin{itemize}
  \item Las cajas \textcolor{azul}{\textbf{azules}} son definiciones, las \textcolor{verde}{\textbf{verdes}} los resultados
  que hay que saber, las \textcolor{naranja}{\textbf{naranjas}} la intuición física y las \textcolor{violeta}{\textbf{violetas}} ejemplos numéricos.
  \item Las cajas \textcolor{rojo}{\textbf{rojas}} (\emph{Ojo}) marcan errores o imprecisiones de las diapositivas de clase,
  o confusiones típicas. La referencia \diapo{n} indica el número de diapositiva.
  \item Convención de signos: corriente del diodo $I$ positiva de $p$ a $n$ por dentro del dispositivo
  (sentido directo), tensión $V_A$ positiva en el lado $p$. La luz agrega una corriente \emph{negativa} (inversa).
\end{itemize}
\end{tcolorbox}

\tableofcontents
\newpage

% =====================================================================
\section{Luz y semiconductores}
% =====================================================================

Los dispositivos optoelectrónicos (fotoconductores, fotodiodos, celdas solares, LEDs y
láseres) son dispositivos semiconductores diseñados para \emph{absorber} o \emph{emitir}
luz de forma eficiente. Todos se basan en la misma idea: un fotón con energía suficiente
puede llevar un electrón de la banda de valencia a la de conducción (\emph{absorción},
genera un par electrón--hueco), y un electrón que cae de la banda de conducción a la de
valencia puede liberar esa energía como un fotón (\emph{emisión}).

\subsection{El fotón: energía y longitud de onda}

La luz está formada por fotones de energía
\begin{equation}
  \Eph = h\nu = \frac{hc}{\lambda},
\end{equation}
con $h=\SI{6.626e-34}{J.s}=\SI{4.136e-15}{\eV.s}$ y $c=\SI{3e8}{m/s}$. Como
$hc = \SI{1.24}{\eV.\micro m}$, en unidades prácticas:

\begin{clave}[Relación energía -- longitud de onda \enhoja]
\[
  \lambda\,[\si{\micro m}] = \frac{1.24}{\Eph\,[\si{\eV}]}
  \qquad\Longrightarrow\qquad
  \lG\,[\si{\micro m}] = \frac{1.24}{\EG\,[\si{\eV}]}
\]
$\lG$ es la \textbf{longitud de onda de corte} del semiconductor: la longitud de onda
cuyo fotón tiene exactamente la energía del \emph{gap}. Fotones con $\lambda<\lG$
($\Eph>\EG$) pueden ser absorbidos; fotones con $\lambda>\lG$ ($\Eph<\EG$) atraviesan el
material: el semiconductor es \textbf{transparente} para ellos.
\end{clave}

\begin{figure}[H]
\centering
\begin{tikzpicture}[x=7.2cm,y=1cm]
  % banda visible con colores reales (modelo "wave" de xcolor, en nm)
  \foreach \l in {380,384,...,748}{
    \definecolor{tmpc}{wave}{\l}
    \fill[tmpc] ({(\l/1000-0.2)},0) rectangle ({(\l/1000-0.2)+0.0042},0.45);
  }
  \fill[violet!25] (0,0) rectangle (0.18,0.45);
  \fill[red!15] (0.552,0) rectangle (1.8,0.45);
  \node[font=\small] at (0.09,0.225) {UV};
  \node[font=\small] at (1.2,0.225) {infrarrojo (IR)};
  \draw[flecha] (0,0) -- (1.85,0) node[right] {$\lambda$ (\si{\micro m})};
  \foreach \x in {0.2,0.4,...,2.0}{
    \draw ({\x-0.2},0) -- ({\x-0.2},-0.1) node[below,font=\scriptsize] {\pgfmathprintnumber[fixed,precision=1]{\x}};
  }
  \node[font=\scriptsize] at (0.37,-0.55) {visible};
  % materiales: lambda_G
  \foreach \lam/\nom/\h in {0.365/GaN/1.0, 0.413/SiC/1.6, 0.46/ZnSe/2.2, 0.549/GaP/1.0,
                             0.873/GaAs/1.0, 0.919/InP/1.6, 1.107/Si/1.0, 1.65/InGaAs/1.0, 1.88/Ge/1.6}{
    \draw[gris,thin] ({\lam-0.2},0.45) -- ({\lam-0.2},\h);
    \node[above,font=\scriptsize,inner sep=1pt] at ({\lam-0.2},\h) {\nom};
  }
  % ventanas de fibra
  \foreach \lam in {1.31,1.55}{
    \draw[azul,densely dashed] ({\lam-0.2},0) -- ({\lam-0.2},-0.9);
  }
  \node[azul,font=\scriptsize,align=center] at (1.23,-1.15) {ventanas de fibra óptica\\1.31 y \SI{1.55}{\micro m}};
  \node[font=\scriptsize,align=left,anchor=west] at (0.95,2.1) {marcas: $\lG=1.24/\EG$ de cada material};
\end{tikzpicture}
\caption{Espectro óptico y longitud de onda de corte $\lG$ de semiconductores comunes (300\,K).
El rango visible va de \SIrange{0.38}{0.75}{\micro m}, es decir de \SIrange{3.3}{1.65}{\eV}.}
\label{fig:espectro}
\end{figure}

\begin{table}[H]
\centering\small
\caption{Gap de algunos semiconductores a 300\,K. La columna ``tipo'' es fundamental: decide si
el material sirve para emitir luz (sección~\ref{sec:led}).}
\label{tab:gaps}
\begin{tabular}{lcccl}
\toprule
Material & $\EG$ (\si{\eV}) & $\lG$ (\si{\micro m}) & Tipo de gap & Uso típico \\
\midrule
Ge                       & 0.66 & 1.88 & indirecto & detectores IR \\
In$_{0.53}$Ga$_{0.47}$As & 0.75 & 1.65 & directo   & detectores para fibra óptica \\
Si                       & 1.12 & 1.11 & indirecto & fotodiodos, celdas solares \\
InP                      & 1.35 & 0.92 & directo   & sustrato, ventanas, láseres \\
GaAs                     & 1.42 & 0.87 & directo   & LEDs IR, láseres, celdas III-V \\
GaP                      & 2.26 & 0.55 & indirecto & LEDs (dopado con N) \\
SiC (6H)                 & 3.0  & 0.41 & indirecto & LEDs azules (antiguos) \\
GaN                      & 3.4  & 0.36 & directo   & LEDs azules/blancos (con InGaN) \\
\bottomrule
\end{tabular}
\end{table}

\begin{ojo}[gap del GaP \diapo{6}]
La tabla de la diapositiva pinta al GaP (\SI{2.3}{\eV}) de verde como si fuera un emisor
natural de verde. El GaP es de \textbf{gap indirecto}: sólo emite con eficiencia razonable
gracias a impurezas \emph{isoelectrónicas} (N, o el par Zn--O) que se agregan a propósito
(ver sección~\ref{sec:led}). Los valores de la diapo están redondeados:
Ge \SI{0.66}{\eV}, Si \SI{1.12}{\eV}, GaAs \SI{1.42}{\eV}.
\end{ojo}

\subsection{Absorción en semiconductores de gap directo e indirecto}

En la absorción hay que conservar la \textbf{energía} y el \textbf{momento cristalino}
($\hbar k$). El momento de un fotón es despreciable: $k_{\text{fotón}}=2\pi/\lambda\sim
\SI{e5}{cm^{-1}}$, mientras que la zona de Brillouin se extiende hasta $k\sim\pi/a\sim
\SI{e8}{cm^{-1}}$. Por eso, en un diagrama $E$--$k$, \emph{la transición debida sólo al
fotón es vertical}.

\begin{figure}[H]
\centering
\begin{tikzpicture}[scale=0.95]
  % ---------- (a) gap directo ----------
  \begin{scope}
    \draw[flecha] (-2,0) -- (2.2,0) node[right] {$k$};
    \draw[flecha] (0,-1.8) -- (0,3.2) node[above] {$E$};
    \draw[bandac,domain=-1.5:1.5,smooth] plot (\x,{1.4+0.6*\x*\x});
    \draw[bandav,domain=-1.6:1.6,smooth] plot (\x,{-0.15-0.45*\x*\x});
    \node[colEc,right] at (1.5,2.75) {BC};
    \node[colEv,right] at (1.6,-1.3) {BV};
    \node[elec] (e) at (0,1.4) {};
    \node[hueco] (h) at (0,-0.15) {};
    \draw[flecha,colFoton,line width=1pt] (0,-0.05) -- (0,1.28);
    \draw[foton] (-1.7,0.9) -- (-0.15,0.65);
    \node[colFoton,font=\small] at (-1.3,1.3) {$h\nu$};
    \draw[<->,gris] (0.35,-0.15) -- node[right,font=\small] {$\EG$} (0.35,1.4);
    \node[font=\small,align=center] at (0,-2.4) {(a) gap \textbf{directo}\\(GaAs, InP, GaN)};
  \end{scope}
  % ---------- (b) gap indirecto ----------
  \begin{scope}[xshift=6.3cm]
    \draw[flecha] (-2,0) -- (2.4,0) node[right] {$k$};
    \draw[flecha] (0,-1.8) -- (0,3.2) node[above] {$E$};
    \draw[bandac,domain=-1.6:2.3,smooth,samples=60] plot (\x,{1.1+0.75*(\x-1.2)*(\x-1.2)/(1+0.35*(\x-1.2)*(\x-1.2))});
    \draw[bandav,domain=-1.6:1.6,smooth] plot (\x,{-0.15-0.45*\x*\x});
    \node[colEc,right] at (2.3,1.6) {BC};
    \node[colEv,right] at (1.6,-1.3) {BV};
    \node[hueco] at (0,-0.15) {};
    \node[elec] at (1.2,1.1) {};
    % foton vertical hasta estado virtual
    \draw[flecha,colFoton,line width=1pt] (0,-0.05) -- (0,0.95);
    \draw[gris,dotted] (0,1.05) -- (1.1,1.05);
    \draw[flecha,violeta,line width=0.9pt] (0.05,1.05) -- (1.08,1.1);
    \node[violeta,font=\small] at (0.75,0.65) {fonón};
    \draw[<->,gris] (1.75,-0.15) -- node[right,font=\small] {$\EG$} (1.75,1.1);
    \draw[gris,dotted] (0,-0.15) -- (1.8,-0.15);
    \node[font=\small,align=center] at (0.2,-2.4) {(b) gap \textbf{indirecto}\\(Si, Ge, GaP)};
  \end{scope}
\end{tikzpicture}
\caption{Absorción de un fotón en los diagramas $E$--$k$. En (a) el mínimo de la BC está
justo encima del máximo de la BV y basta el fotón. En (b) el electrón tiene que cambiar de
$k$: necesita además \emph{absorber o emitir un fonón} (vibración de la red), que aporta
el momento faltante (y una pequeña energía $E_{ph}\sim\SI{50}{meV}$).}
\label{fig:Ek}
\end{figure}

\begin{intuicion}[por qué importa el tipo de gap]
Un proceso que necesita fotón \emph{y} fonón a la vez es mucho menos probable que uno que
necesita sólo el fotón. Consecuencias:
\begin{itemize}
  \item \textbf{Absorción:} en un material indirecto el coeficiente de absorción $\alpha$
  crece lentamente al pasar $\lG$; hacen falta espesores grandes (en Si, una celda solar
  necesita $\sim\SI{100}{\micro m}$ o más; en GaAs alcanzan unos pocos \si{\micro m}).
  \item \textbf{Emisión:} es el mismo argumento al revés. En Si la recombinación
  banda a banda radiativa es tan lenta que gana la recombinación por trampas, que
  libera la energía como calor. Por eso \emph{no hay LEDs de silicio}.
\end{itemize}
Con absorción de fonón se pueden absorber fotones con $h\nu=\EG-E_{ph}$; con emisión de
fonón, $h\nu=\EG+E_{ph}$ (las tres variantes del diagrama de la \diapo{3}).
\end{intuicion}

\subsection{Coeficiente de absorción y ley de Beer--Lambert}

Supongamos un haz de luz monocromática que entra a un semiconductor por $x=0$. Cada
unidad de longitud recorrida absorbe una fracción fija $\alpha$ de los fotones
que quedan:
\[
  \frac{\dd\Phi}{\dd x} = -\alpha\,\Phi
  \qquad\Longrightarrow\qquad
  \Phi(x) = \Phi_0\,e^{-\alpha x},\qquad I(x)=I_0\,e^{-\alpha x},
\]
donde $\Phi$ es el flujo de fotones (fotones/\si{cm^2.s}) e $I$ la intensidad
(\si{W/cm^2}). $\alpha$ [\si{cm^{-1}}] es el \textbf{coeficiente de absorción}, y
$1/\alpha$ es la \textbf{profundidad de penetración}: la distancia en que la luz cae a
$1/e\approx\SI{37}{\percent}$. A una profundidad de $3/\alpha$ ya se absorbió el
\SI{95}{\percent}.

\begin{figure}[H]
\centering
\begin{tikzpicture}
\begin{semilogyaxis}[
  width=0.62\textwidth, height=6.4cm,
  xmin=0.3, xmax=1.9, ymin=1, ymax=2e6,
  xlabel={$\lambda$ (\si{\micro m})}, ylabel={$\alpha$ (\si{cm^{-1}})},
  grid=both, minor grid style={gray!15}, major grid style={gray!30},
  legend pos=south west, legend style={font=\small},
  every axis plot/.append style={line width=1.1pt,smooth},
  axis y line*=left,
]
\addplot[azul] coordinates {(0.30,1.7e6)(0.35,1.0e6)(0.40,9.5e4)(0.45,2.6e4)(0.50,1.1e4)(0.55,6.4e3)(0.60,4.1e3)
  (0.65,2.8e3)(0.70,1.9e3)(0.75,1.3e3)(0.80,850)(0.85,535)(0.90,306)(0.95,157)(1.00,64)(1.05,16)(1.10,3.5)(1.12,1.2)};
\addlegendentry{Si}
\addplot[rojo] coordinates {(0.30,1.5e6)(0.40,7e5)(0.50,1.2e5)(0.60,5e4)(0.70,2.8e4)(0.80,1.3e4)(0.85,8e3)
  (0.865,3e3)(0.875,4e2)(0.885,30)(0.89,1)};
\addlegendentry{GaAs}
\addplot[verde] coordinates {(0.40,5e5)(0.50,2e5)(0.60,8e4)(0.70,5e4)(0.80,3e4)(0.90,2e4)(0.915,6e3)(0.925,5e2)(0.935,10)(0.94,1)};
\addlegendentry{InP}
\addplot[naranja] coordinates {(0.40,6e5)(0.60,2e5)(0.80,6e4)(1.00,2e4)(1.20,1e4)(1.40,7e3)(1.50,4e3)(1.55,7e2)
  (1.60,1.5e2)(1.70,30)(1.80,5)(1.86,1)};
\addlegendentry{Ge}
\end{semilogyaxis}
\begin{semilogyaxis}[
  width=0.62\textwidth, height=6.4cm,
  xmin=0.3, xmax=1.9, ymin=1, ymax=2e6,
  axis x line=none, axis y line*=right,
  ytick={1e-4,1e-3,1e-2,1e-1,1,10,100,1000,10000,100000,1000000},
  yticklabels={},
  extra y ticks={1e2,1e3,1e4,1e5},
  extra y tick labels={\SI{100}{\micro m},\SI{10}{\micro m},\SI{1}{\micro m},\SI{0.1}{\micro m}},
  extra y tick style={tick label style={font=\scriptsize}},
  ylabel={$1/\alpha$}, ylabel style={font=\small},
]
\end{semilogyaxis}
\end{tikzpicture}
\caption{Coeficiente de absorción en función de $\lambda$ (valores aproximados a 300\,K,
adaptado de Pierret, fig.~3.20). A la derecha, la profundidad de penetración $1/\alpha$
correspondiente. Notar el corte \emph{abrupto} de los materiales directos (GaAs, InP) en
$\lG$ y el corte \emph{gradual} del Si (indirecto). El Ge tiene un gap indirecto de
\SI{0.66}{\eV} y uno directo de \SI{0.80}{\eV} ($\approx\SI{1.55}{\micro m}$): de ahí el
escalón.}
\label{fig:alfa}
\end{figure}

\begin{table}[H]
\centering\small
\caption{Profundidad de penetración $1/\alpha$ en silicio. La luz azul se absorbe en la
superficie; la infrarroja cercana penetra cientos de micrones.}
\label{tab:penetracion}
\begin{tabular}{lccccccc}
\toprule
$\lambda$ (\si{\micro m}) & 0.40 & 0.50 & 0.60 & 0.70 & 0.80 & 0.90 & 1.00 \\
\midrule
$\alpha$ (\si{cm^{-1}}) & \num{9.5e4} & \num{1.1e4} & \num{4.1e3} & \num{1.9e3} & 850 & 306 & 64 \\
$1/\alpha$ (\si{\micro m}) & 0.1 & 0.9 & 2.4 & 5.3 & 12 & 33 & 156 \\
\bottomrule
\end{tabular}
\end{table}

\subsection{Fotogeneración}
\label{sec:fotogeneracion}

Cada fotón absorbido (con $\Eph>\EG$) genera un par electrón--hueco. La tasa de
fotogeneración por unidad de volumen es entonces el número de fotones absorbidos por
unidad de volumen y de tiempo, $-\dd\Phi/\dd x$:

\begin{clave}[Tasa de fotogeneración \enhoja]
\[
  G_L(x,\lambda) = \left.\frac{\partial n}{\partial t}\right|_{\text{luz}}
  = \left.\frac{\partial p}{\partial t}\right|_{\text{luz}}
  = \alpha\,\Phi_0\,e^{-\alpha x} \equiv G_{L0}\,e^{-\alpha x}
  \qquad [\si{cm^{-3}.s^{-1}}]
\]
con $\Phi_0 = \dfrac{(1-R)\,P}{h\nu\,A}$ el flujo de fotones que \emph{entra} al
semiconductor (potencia $P$ sobre área $A$, reflectividad $R$). Electrones y huecos se
generan \textbf{de a pares}: $\Delta n=\Delta p$ en el instante de la generación.
\end{clave}

\begin{figure}[H]
\centering
\begin{tikzpicture}
  % bloque semiconductor
  \fill[azulclaro] (0,0) rectangle (6,1.4);
  \draw (0,0) rectangle (6,1.4);
  \node[font=\small,anchor=south east] at (5.95,0.05) {semiconductor};
  % sombreado de intensidad decreciente
  \foreach \i in {0,...,29}{
    \pgfmathsetmacro{\op}{0.85*exp(-\i/9)}
    \fill[naranja,opacity=\op] ({\i*0.2},0) rectangle ({\i*0.2+0.2},1.4);
  }
  \foreach \y in {0.3,0.7,1.1}{
    \draw[foton] (-1.6,\y) -- (-0.05,\y);
  }
  \node[colFoton,above] at (-0.9,1.4) {$h\nu$, $\lambda$};
  % grafico debajo
  \begin{scope}[yshift=-3.2cm]
    \draw[flecha] (0,0) -- (6.4,0) node[right] {$x$};
    \draw[flecha] (0,0) -- (0,2.4) node[above] {$G_L(x)$};
    \draw[naranja,line width=1.1pt,domain=0:6,smooth,samples=60] plot (\x,{2*exp(-\x/1.8)});
    \draw[gris,dashed] (1.8,0) node[below,black] {$1/\alpha$} -- (1.8,{2*exp(-1)});
    \draw[gris,dashed] (0,{2*exp(-1)}) node[left,black] {$G_{L0}/e$} -- (1.8,{2*exp(-1)});
    \node[left] at (0,2) {$G_{L0}$};
    \node[font=\small,align=left,anchor=west] at (3.0,1.4) {$G_L=G_{L0}\,e^{-\alpha x}$\\$\alpha$ grande: todo en la superficie\\$\alpha$ chico: casi uniforme};
  \end{scope}
\end{tikzpicture}
\caption{La luz se atenúa exponencialmente al penetrar y con ella la tasa de generación de
pares \diapo{4}. Si el dispositivo es mucho más delgado que $1/\alpha$, se puede aproximar
$G_L\approx$ constante, que es la hipótesis que usa Pierret en las deducciones de
fotodiodos.}
\label{fig:fotogeneracion}
\end{figure}

\paragraph{Reflexión en la superficie.} Al pasar del aire ($n_1=1$) a un semiconductor de
índice $n$ con incidencia normal se refleja una fracción
\[
  R = \left(\frac{n-1}{n+1}\right)^2 .
\]
Para Si ($n\approx3.5$) da $R\approx\SI{31}{\percent}$: ¡casi un tercio de la luz se pierde
antes de entrar! Por eso todos los fotodetectores y celdas llevan \textbf{capa
antirreflectante}: una película de índice $n_{AR}\approx\sqrt{n}$ y espesor
$\lambda/(4n_{AR})$ hace que las ondas reflejadas en sus dos caras se cancelen. Las celdas
solares además \emph{texturizan} la superficie (pirámides) para que la luz rebotada
vuelva a incidir sobre el semiconductor.

\begin{ejemplo}[¿qué material para qué luz?]
\textbf{a)} ¿Puede un fotodiodo de Si detectar la luz de un láser de fibra de
\SI{1.55}{\micro m}? \quad
$\Eph = 1.24/1.55 = \SI{0.80}{\eV} < \EG(\text{Si}) = \SI{1.12}{\eV}$. \textbf{No}: el Si
es transparente a esa luz. Sirven Ge ($\lG=\SI{1.88}{\micro m}$) o
In$_{0.53}$Ga$_{0.47}$As ($\lG=\SI{1.65}{\micro m}$).

\textbf{b)} Luz roja de \SI{0.65}{\micro m} incide sobre Si. ¿Qué espesor absorbe el
\SI{90}{\percent} de lo que entra? \quad Interpolando la tabla~\ref{tab:penetracion},
$\alpha\approx\SI{2.8e3}{cm^{-1}}$. Se pide $1-e^{-\alpha d}=0.9\Rightarrow
d=\ln 10/\alpha = 2.3/\SI{2.8e3}{cm^{-1}}\approx\SI{8}{\micro m}$.
\end{ejemplo}

\begin{ojo}[\diapo{5}]
La frase ``el fotodiodo PN tiene una respuesta en frecuencia limitada a aprox.\
\SI{10}{MHz}'' aparece en la diapositiva del coeficiente de absorción pero no tiene que ver
con $\alpha$ directamente: se explica en la sección~\ref{sec:pn-velocidad} (es la lentitud
de la difusión de los portadores generados fuera de la zona de carga espacial).
\end{ojo}

% =====================================================================
\section{Recombinación y generación}
\label{sec:rg}
% =====================================================================

En equilibrio térmico la generación y la recombinación de pares se compensan
($np=n_i^2$). Cuando la luz (o la inyección de una juntura) saca al semiconductor del
equilibrio, los mecanismos de recombinación--generación (R--G) son los que intentan
devolverlo. Conocerlos explica, por ejemplo, por qué el Si no emite luz y por qué la
fotocorriente depende de la vida media $\tau$.

\subsection{Mecanismos}

\begin{figure}[H]
\centering
\begin{tikzpicture}[scale=0.92,
  ec/.style={bandac},ev/.style={bandav}]
  % helper: dibuja bandas
  \newcommand{\bandas}[1]{%
    \draw[ec] (#1,2) -- ++(2.6,0) node[right,font=\scriptsize] {$\Ec$};
    \draw[ev] (#1,0) -- ++(2.6,0) node[right,font=\scriptsize] {$\Ev$};}
  % ======= fila superior: recombinacion =======
  \node[font=\bfseries\small,anchor=west] at (-1.6,3.2) {Recombinación};
  % (a) banda a banda
  \bandas{0}
  \node[elec] at (0.9,2) {}; \node[hueco] at (0.9,0) {};
  \draw[flecha] (0.9,1.9) -- (0.9,0.12);
  \draw[foton] (1.05,1) -- (2.4,1);
  \node[colFoton,font=\scriptsize] at (1.75,0.7) {fotón};
  \node[font=\scriptsize,align=center] at (1.3,-0.6) {(a) banda a banda\\(directa, \textbf{radiativa})};
  % (b) centro R-G
  \bandas{4}
  \draw[fermi] (4,1) -- (6.6,1) node[right,font=\scriptsize,black] {$E_T$};
  \node[elec] at (4.6,2) {}; \node[hueco] at (5.7,0) {};
  \draw[flecha] (4.6,1.9) -- (4.6,1.05);
  \node[elec] at (4.6,1) {};
  \draw[flecha,dashed] (4.75,1) -- (5.6,1);
  \draw[flecha] (5.7,0.95) -- (5.7,0.12);
  \draw[flecha,decorate,decoration={snake,amplitude=1pt,segment length=4pt},gris] (4.75,1.5) -- (5.35,1.5);
  \node[gris,font=\scriptsize,anchor=west] at (5.35,1.5) {calor};
  \node[font=\scriptsize,align=center] at (5.3,-0.6) {(b) por centros R--G\\(trampas, SRH)};
  % (c) Auger
  \bandas{8}
  \node[elec] (e1) at (8.7,2) {}; \node[elec] at (9.3,2) {};
  \node[hueco] at (8.7,0) {};
  \draw[flecha] (8.7,1.9) -- (8.7,0.12);
  \draw[flecha] (9.3,2.05) -- (9.3,2.9) node[above,font=\scriptsize] {$e^-$ ``caliente''};
  \foreach \yy in {2.75,2.55,2.35,2.15}{\draw[gris] (9.55,\yy) -- ++(0.25,0);}
  \node[font=\scriptsize,align=center] at (9.3,-0.6) {(c) Auger\\(3 portadores)};
  % ======= fila inferior: generacion =======
  \begin{scope}[yshift=-4.4cm]
  \node[font=\bfseries\small,anchor=west] at (-1.6,3.2) {Generación};
  \bandas{0}
  \node[elec] at (0.9,2) {}; \node[hueco] at (0.9,0) {};
  \draw[flecha] (0.9,0.12) -- (0.9,1.9);
  \draw[foton] (2.4,1) -- (1.05,1);
  \node[colFoton,font=\scriptsize,align=center] at (1.85,1.55) {luz o calor};
  \node[font=\scriptsize,align=center] at (1.3,-0.6) {(d) banda a banda\\(\textbf{fotogeneración})};
  \bandas{4}
  \draw[fermi] (4,1) -- (6.6,1) node[right,font=\scriptsize,black] {$E_T$};
  \node[elec] at (4.6,2) {}; \node[hueco] at (5.7,0) {};
  \draw[flecha] (4.6,1.05) -- (4.6,1.9);
  \draw[flecha] (5.7,0.12) -- (5.7,0.95);
  \node[font=\scriptsize,align=center] at (5.3,-0.6) {(e) por centros R--G\\(térmica)};
  \bandas{8}
  \draw[flecha] (8.3,2.9) node[elec] {} -- (9.0,2.1);
  \node[elec] at (9.2,2) {}; \node[elec] at (9.8,2) {}; \node[hueco] at (9.8,0) {};
  \draw[flecha] (9.8,0.12) -- (9.8,1.9);
  \node[font=\scriptsize,align=center] at (9.3,-0.6) {(f) ionización por impacto\\(avalancha)};
  \end{scope}
\end{tikzpicture}
\caption{Procesos de recombinación (arriba) y generación (abajo) \diapo{2}. Electrones:
círculos llenos; huecos: círculos vacíos. Cada proceso de generación es el inverso del
de recombinación que tiene encima.}
\label{fig:rg}
\end{figure}

\begin{description}[style=nextline,leftmargin=1.2em]
  \item[Banda a banda (radiativa)] El electrón cae directamente de la BC a la BV y emite un
  fotón con $h\nu\approx\EG$. Es el mecanismo dominante en semiconductores de \emph{gap
  directo} y es la base del LED y del láser. Su inverso es la \textbf{fotogeneración}.
  \item[Por centros R--G (Shockley--Read--Hall)] Impurezas o defectos introducen niveles
  $E_T$ cerca del medio del gap. El electrón cae en dos escalones y la energía se entrega a
  la red como vibraciones (\textbf{calor}). Domina en Si y Ge (gap indirecto). Es el
  mecanismo que fija la vida media $\tau$ en dispositivos de Si.
  \item[Auger] La energía de la recombinación la recibe un \emph{tercer} portador, que
  queda ``caliente'' y la pierde luego por choques. Importa con dopajes o
  inyecciones muy altas (limita la eficiencia de LEDs a alta corriente y de celdas de Si
  de alta calidad).
  \item[Ionización por impacto] Un portador muy acelerado por un campo intenso choca con
  la red y arranca un electrón de la BV. Es la base de la ruptura por avalancha y del
  fotodiodo de avalancha (sección~\ref{sec:apd}).
\end{description}

\subsection{Vida media y ecuación de continuidad}

Para \emph{baja inyección} (perturbación pequeña frente a los mayoritarios), la
recombinación neta de minoritarios es proporcional al exceso. Por ejemplo, para
electrones en un material tipo p:
\[
  \left.\frac{\partial \Delta n}{\partial t}\right|_{\text{R--G}} = -\frac{\Delta n}{\tau_n},
\]
donde $\tau_n$ es la \textbf{vida media} de los minoritarios: el tiempo promedio que un
electrón en exceso sobrevive antes de recombinarse. Juntando deriva, difusión,
recombinación y luz se obtiene la ecuación de difusión de minoritarios:

\begin{clave}[Ecuación de difusión de minoritarios con luz \enhoja]
\[
  \frac{\partial \Delta n_p}{\partial t}
  = D_n\frac{\partial^2 \Delta n_p}{\partial x^2}
  - \frac{\Delta n_p}{\tau_n}
  + G_L
\]
(análoga para huecos en material n con $D_p$, $\tau_p$).
La \textbf{longitud de difusión} $L_n=\sqrt{D_n\tau_n}$ es la distancia media que difunde
un minoritario antes de recombinarse.
\end{clave}

\paragraph{Caso uniforme y estacionario.} Si $G_L$ es uniforme y no hay gradientes,
$0 = -\Delta n/\tau_n + G_L$, es decir
\begin{equation}
  \boxed{\Delta n = \Delta p = G_L\,\tau_n}
  \label{eq:dn-estacionario}
\end{equation}
El exceso de portadores es ``lo que se genera por segundo'' multiplicado por ``cuánto
dura cada uno''. Este resultado es la base del fotoconductor.

\paragraph{Transitorio.} Si se enciende la luz en $t=0$ y se apaga en $t=t_1$:
\[
  \Delta n(t) = G_L\tau_n\left(1-e^{-t/\tau_n}\right)\quad(0<t<t_1),
  \qquad
  \Delta n(t) = \Delta n(t_1)\,e^{-(t-t_1)/\tau_n}\quad(t>t_1).
\]

\begin{figure}[H]
\centering
\begin{tikzpicture}
\begin{axis}[
  width=0.7\textwidth,height=4.6cm,
  xmin=0,xmax=12,ymin=0,ymax=1.35,
  xlabel={$t/\tau_n$},ylabel={$\Delta n/(G_L\tau_n)$},
  xtick={0,2,4,6,8,10,12},ytick={0,0.5,1},
  axis lines=left,
]
\fill[naranjaclaro] (axis cs:1,0) rectangle (axis cs:6,1.35);
\node[naranja,font=\small] at (axis cs:3.5,1.22) {luz encendida};
\addplot[azul,line width=1.2pt,domain=0:1] {0};
\addplot[azul,line width=1.2pt,domain=1:6,samples=60] {1-exp(-(x-1))};
\addplot[azul,line width=1.2pt,domain=6:12,samples=60] {(1-exp(-5))*exp(-(x-6))};
\draw[gris,dashed] (axis cs:0,1) -- (axis cs:12,1);
\draw[gris,dashed] (axis cs:2,0) -- (axis cs:2,{1-exp(-1)});
\node[gris,font=\scriptsize,anchor=west] at (axis cs:2.1,0.45) {$\tau_n$: \SI{63}{\percent}};
\end{axis}
\end{tikzpicture}
\caption{Respuesta del exceso de portadores a un pulso de luz. La vida media fija tanto el
nivel de señal ($G_L\tau_n$) como la velocidad de respuesta: hay un compromiso
\emph{sensibilidad--velocidad} que reaparece en el fotoconductor.}
\label{fig:transitorio}
\end{figure}

% =====================================================================
\section{Fotoconductores}
\label{sec:fotoconductor}
% =====================================================================

El fotoconductor es el fotodetector más simple: una pastilla de semiconductor
homogéneo (tipo n o p, sin juntura) con dos contactos óhmicos. La luz genera pares,
aumenta la conductividad y, a tensión fija, aumenta la corriente. Ejemplos típicos son las
LDR (\emph{light dependent resistors}) de CdS o CdSe para luz visible, y las de PbS o
PbSe para infrarrojo.

\begin{figure}[H]
\centering
\begin{circuitikz}[scale=0.95,transform shape]
  % pastilla semiconductora (vista en perspectiva)
  \fill[azulclaro] (3,1) rectangle (4.2,4);
  \fill[azul!35] (4.2,1) -- (4.6,1.4) -- (4.6,4.4) -- (4.2,4) -- cycle;
  \fill[gris!60] (3,4) -- (4.2,4) -- (4.6,4.4) -- (3.4,4.4) -- cycle;
  \draw (3,1) rectangle (4.2,4) (4.2,1) -- (4.6,1.4) -- (4.6,4.4) -- (4.2,4) (3,4) -- (3.4,4.4) -- (4.6,4.4);
  \fill[gris!60] (3,1) rectangle (4.2,0.85);
  \draw (3,1) rectangle (4.2,0.85);
  \node[font=\small,anchor=south west] at (4.6,4.3) {área $A$};
  \draw[<->] (2.75,1) -- node[left] {$L$} (2.75,4);
  \foreach \y in {1.8,2.5,3.2}{\draw[foton] (6.1,\y+0.3) -- (4.75,\y);}
  \node[colFoton] at (6.5,3.6) {luz};
  % circuito
  \draw (3.8,4.2) -- (3.8,5.2) to[short,i_=$I$] (0,5.2) to[battery1,l_=$V$] (0,-1.4) node[ground]{};
  \draw (3.6,0.85) -- (3.6,0.2) to[R,l_=$R_L$,*-] (3.6,-1.4) node[ground]{};
  \draw (3.6,0.2) -- (5,0.2) to[C,l=$C$,-o] (7,0.2) node[right] {$v_o$};
\end{circuitikz}
\caption{Fotoconductor polarizado con una fuente $V$ y una resistencia de carga $R_L$
\diapo{44}. El capacitor deja pasar a la salida sólo la parte variable de la señal
(la luz modulada), bloqueando la corriente de oscuridad.}
\label{fig:fotoconductor}
\end{figure}

\subsection{Cambio de conductividad}

Sin luz la conductividad es $\sigma_0=q(\mu_n n_0+\mu_p p_0)$. Con luz aparecen excesos
$\Delta n=\Delta p$ (se generan de a pares):
\[
  \sigma = q\left[\mu_n(n_0+\Delta n)+\mu_p(p_0+\Delta p)\right]
  \quad\Longrightarrow\quad
  \Delta\sigma = q(\mu_n+\mu_p)\,\Delta n .
\]
En estado estacionario y con generación uniforme, por \eqref{eq:dn-estacionario},
$\Delta n=G_L\tau$, donde $\tau$ es la vida media de los minoritarios
(supondremos material tipo p, así que $\tau=\tau_n$). La corriente adicional con una
tensión $V$ aplicada sobre la longitud $L$ (campo $\mathcal{E}=V/L$) es
\begin{equation}
  \Delta I = \Delta\sigma\,\mathcal{E}\,A = q(\mu_n+\mu_p)\,G_L\tau_n\,\frac{A\,V}{L}.
  \label{eq:dI-fotoconductor}
\end{equation}

\begin{ojo}[resistencia del fotoconductor \diapo{44}]
La diapositiva escribe $\Delta R = L/(\Delta\sigma A)$. Eso \textbf{no} es la variación de
resistencia: $L/(\Delta\sigma A)$ es la resistencia de un conductor ``imaginario'' que sólo
tuviera la conductividad agregada por la luz. Lo que es aditivo es la
\emph{conductancia}, porque ambas conducciones están en paralelo:
\[
  G_0=\frac{\sigma_0 A}{L},\qquad G=\frac{(\sigma_0+\Delta\sigma)A}{L}
  \;\Rightarrow\; \Delta G = \frac{\Delta\sigma\,A}{L},
  \qquad
  R=\frac{L}{(\sigma_0+\Delta\sigma)A}=\frac{R_0}{1+\Delta\sigma/\sigma_0}.
\]
La resistencia \emph{disminuye} con la luz; la variación real es
$\Delta R = R-R_0 = -R_0\,\dfrac{\Delta\sigma}{\sigma_0+\Delta\sigma}$.
\end{ojo}

\subsection{Eficiencia cuántica y ganancia}

\begin{definicion}[eficiencia cuántica $\eta$ \enhoja]
\[
  \eta = \frac{\text{pares e--h generados por segundo}}{\text{fotones incidentes por segundo}}
  = \frac{G_L\,A\,L}{\Popt/h\nu}
  \qquad (0\le\eta\le 1)
\]
$G_L A L$ es la generación total en el volumen $A\,L$ del dispositivo y $\Popt/h\nu$ el
número de fotones que llegan por segundo (potencia óptica sobre energía por fotón).
\end{definicion}

Si cada par generado aportara \emph{una} carga $q$ al circuito, la corriente sería
\[
  I_n \equiv q\,G_L A L = q\,\eta\,\frac{\Popt}{h\nu}
  \qquad\text{(``corriente primaria'')}.
\]
Pero el fotoconductor entrega más que eso:

\begin{definicion}[ganancia óptica (fotoconductiva) \enhoja]
\[
  G_{opt} = \frac{\Delta I}{I_n}
  = \frac{\text{corriente que entrega el dispositivo}}{\text{carga generada por la luz por unidad de tiempo}}
\]
\end{definicion}

Usando \eqref{eq:dI-fotoconductor}:
\[
  G_{opt} = \frac{q(\mu_n+\mu_p)G_L\tau_n\,AV/L}{q\,G_L A L}
  = \frac{(\mu_n+\mu_p)\,\tau_n V}{L^2}.
\]
Definiendo el \textbf{tiempo de tránsito} de un electrón entre contactos,
$t_{tr,n} = L/v_n = L/(\mu_n\mathcal{E}) = L^2/(\mu_n V)$:

\begin{clave}[Ganancia del fotoconductor \enhoja]
\[
  G_{opt} = \frac{\tau_n}{t_{tr,n}}\left(1+\frac{\mu_p}{\mu_n}\right)
  \;\approx\; \frac{\tau_n}{t_{tr,n}} = \frac{\tau_n\,\mu_n\,\mathcal{E}}{L}
  \qquad(\text{si }\mu_n\gg\mu_p)
\]
La ganancia es el cociente entre cuánto vive un portador y cuánto tarda en cruzar el
dispositivo.
\end{clave}

\begin{ojo}[\diapo{46}]
La diapositiva llega a $G_{opt}=\tau_n\mu_n\mathcal{E}/L$ partiendo de una expresión que
tiene $(\mu_n+\mu_p)$. Al reemplazar se ``pierde'' el término de huecos: el resultado
exacto (con las mismas hipótesis) es $G_{opt}=\tau_n(\mu_n+\mu_p)\mathcal{E}/L$. La forma
de la diapo es una aproximación válida si $\mu_n\gg\mu_p$ o si los huecos quedan atrapados
y no contribuyen. Además, en la expresión de $I_n$ de la diapo~45 sobra el ``$\cdot$'' antes
del igual (errata).
\end{ojo}

\begin{intuicion}[¿cómo puede haber ganancia mayor que 1?]
Un electrón fotogenerado cruza la muestra en $t_{tr}$ y sale por el contacto positivo.
Como la muestra tiene que seguir siendo neutra (su hueco compañero sigue adentro), el
contacto negativo \emph{inyecta otro electrón} de inmediato. Ese electrón vuelve a cruzar,
y así sucesivamente, hasta que el hueco se recombina (en un tiempo $\sim\tau$). Cada par
fotogenerado hace circular $\tau/t_{tr}$ electrones por el circuito.

El precio es la velocidad: la señal tarda $\sim\tau$ en subir y bajar
(figura~\ref{fig:transitorio}). Aumentar $\tau$ aumenta la ganancia pero empeora el ancho
de banda $B\sim1/(2\pi\tau)$; el producto \emph{ganancia $\times$ ancho de banda}
$\sim 1/(2\pi t_{tr})$ queda fijo por la geometría y la tensión. Por eso los
fotoconductores son sensibles pero lentos (típicamente ms en CdS).
\end{intuicion}

\begin{ejemplo}[ganancia de un fotoconductor de Si]
Pastilla de Si tipo p, $L=\SI{1}{mm}$, $V=\SI{10}{V}$, $\mu_n=\SI{1350}{cm^2/V.s}$,
$\mu_p=\SI{480}{cm^2/V.s}$, $\tau_n=\SI{10}{\micro s}$.
\[
  t_{tr,n}=\frac{L^2}{\mu_n V}=\frac{(\SI{0.1}{cm})^2}{\SI{1350}{cm^2/V.s}\cdot\SI{10}{V}}
  =\SI{0.74}{\micro s},
  \qquad
  G_{opt}=\frac{\SI{10}{\micro s}}{\SI{0.74}{\micro s}}\left(1+\frac{480}{1350}\right)\approx 18 .
\]
(Con la fórmula aproximada daría $\approx13.5$.) El tiempo de respuesta es del orden de
$\tau_n=\SI{10}{\micro s}$, o sea un ancho de banda de apenas
$\sim1/(2\pi\tau_n)\approx\SI{16}{kHz}$.
\end{ejemplo}

% =====================================================================
\section{Fotodiodo de juntura PN}
\label{sec:pn}
% =====================================================================

\subsection{Principio de funcionamiento}

Un fotodiodo es una juntura PN diseñada para que la luz llegue a la zona de la juntura.
La clave es el \textbf{campo eléctrico de la zona de carga espacial (ZCE)}: separa los
pares fotogenerados (electrones hacia el lado n, huecos hacia el lado p) y los convierte
en corriente antes de que se recombinen. Como el campo apunta de n a p, la fotocorriente
circula en sentido \textbf{inverso} (de n a p por dentro del diodo).

Según dónde se genere el par, hay tres casos (figura~\ref{fig:pn-colecta}):
\begin{enumerate}
  \item \textbf{Dentro de la ZCE} (ancho $W$): el campo separa el par de inmediato
  (\emph{deriva}, muy rápido). Se colecta todo.
  \item \textbf{En la zona neutra, a menos de una longitud de difusión de la ZCE}: el
  minoritario (electrón en p, hueco en n) \emph{difunde} al azar; con buena probabilidad
  llega al borde de la ZCE y el campo lo barre al otro lado. Es un proceso lento.
  \item \textbf{Lejos de la ZCE} ($\gg L_n$ o $L_p$): el minoritario se recombina antes de
  llegar. No aporta corriente.
\end{enumerate}

\begin{figure}[H]
\centering
\begin{tikzpicture}[scale=1]
  % regiones
  \fill[rojoclaro] (0,0) rectangle (5,1.6);
  \fill[azulclaro] (6.2,0) rectangle (11.2,1.6);
  \fill[gris!15] (5,0) rectangle (6.2,1.6);
  \draw (0,0) rectangle (11.2,1.6);
  \draw[dashed] (5,0) -- (5,1.6) (6.2,0) -- (6.2,1.6);
  \node at (0.5,1.3) {p}; \node at (10.8,1.3) {n};
  % zonas de colecta
  \fill[verde,opacity=0.18] (2.6,0) rectangle (8.6,1.6);
  \draw[<->] (2.6,-0.3) -- node[below,font=\small] {$L_n$} (5,-0.3);
  \draw[<->] (5,-0.3) -- node[below,font=\small] {$W$} (6.2,-0.3);
  \draw[<->] (6.2,-0.3) -- node[below,font=\small] {$L_p$} (8.6,-0.3);
  \node[verde!70!black,font=\small] at (5.6,-1.0) {volumen de colecta $A(L_n+W+L_p)$};
  % campo
  \draw[flecha,line width=1pt] (6.05,0.25) -- (5.15,0.25);
  \node[font=\scriptsize] at (5.6,0.48) {$\vec{\mathcal{E}}$};
  % pares en ZCE
  \node[elec] at (5.4,1.0) {}; \node[hueco] at (5.8,1.0) {};
  \draw[flecha,colElec] (5.55,1.15) -- (6.1,1.15);
  \draw[flecha,colHueco] (5.65,0.85) -- (5.1,0.85);
  % par en p cerca: electron difunde
  \node[elec] (e1) at (3.6,0.9) {};
  \draw[colElec,decorate,decoration={random steps,segment length=3pt,amplitude=2pt},-{Stealth[length=4pt]}] (3.65,0.9) -- (4.95,1.05);
  % par en n cerca: hueco difunde
  \node[hueco] at (7.6,0.7) {};
  \draw[colHueco,decorate,decoration={random steps,segment length=3pt,amplitude=2pt},-{Stealth[length=4pt]}] (7.5,0.7) -- (6.25,0.6);
  % par lejos: recombina
  \node[elec] at (1.0,0.6) {}; \node[hueco] at (1.25,0.6) {};
  \node[font=\scriptsize,align=center] at (1.15,0.25) {se recombina};
  \node[elec] at (10.0,0.65) {}; \node[hueco] at (10.25,0.65) {};
  \node[font=\scriptsize,align=center] at (10.15,0.25) {se recombina};
  % luz
  \foreach \x in {1.2,3.6,5.6,7.6,10}{\draw[foton] (\x,2.9) -- (\x,1.7);}
  \node[colFoton,anchor=west] at (10.2,2.6) {luz};
\end{tikzpicture}
\caption{Colección de portadores en un fotodiodo PN (adaptado de Pierret, fig.~9.2 y
\diapo{8}). Sólo los pares generados en la ZCE o a menos de una longitud de difusión de ella
contribuyen a la corriente.}
\label{fig:pn-colecta}
\end{figure}

\subsection{Característica I--V con iluminación}

\begin{clave}[Ecuación del fotodiodo \enhoja]
\[
  I = I_0\left(e^{V_A/V_T}-1\right) - I_L,
  \qquad
  I_L = q\,A\,G_L\,(L_n + W + L_p)
\]
\textbf{En la hoja de fórmulas} (y en Pierret) está escrita como $I=I_S(e^{V/V_T}-1)+I_L$
con $I_L=-qAG_L(L_n+W+L_p)$: es lo mismo, sólo que ahí $I_L$ ya incluye el signo
negativo. Hay que tener claro cuál convención se está usando al hacer cuentas. La luz \textbf{desplaza la curva del diodo hacia abajo}
una cantidad $I_L$ proporcional a la intensidad de la luz, sin cambiar su forma.
\end{clave}

\paragraph{Deducción (ecuación de difusión con generación).} Tomemos el lado n
($x\ge0$ medido desde el borde de la ZCE) con generación uniforme $G_L$. En estado
estacionario la ecuación de difusión para huecos minoritarios queda
\[
  0 = D_p\frac{\dd^2\Delta p}{\dd x^2} - \frac{\Delta p}{\tau_p} + G_L,
  \qquad
  \Delta p(0) = p_{n0}\left(e^{V_A/V_T}-1\right),
  \quad
  \Delta p(\infty) = G_L\tau_p .
\]
La primera condición es la ley de la juntura; la segunda dice que lejos de la juntura el
exceso es el de una muestra uniformemente iluminada \eqref{eq:dn-estacionario}. La solución es
\[
  \Delta p(x) = G_L\tau_p + \left[p_{n0}\left(e^{V_A/V_T}-1\right)-G_L\tau_p\right]e^{-x/L_p}.
\]
La corriente de huecos en el borde de la ZCE es de difusión:
\[
  J_p(0) = -qD_p\left.\frac{\dd\Delta p}{\dd x}\right|_0
  = \frac{qD_p}{L_p}\,p_{n0}\left(e^{V_A/V_T}-1\right) - q\,G_L\,L_p ,
\]
donde se usó $D_p\tau_p/L_p=L_p$. El primer término es la corriente del diodo ideal; el
segundo, \textbf{la generación dentro de una longitud $L_p$}. El lado p aporta análogamente
$-qG_LL_n$, y la ZCE, donde todo lo generado se colecta, $-qG_LW$. Sumando y
multiplicando por $A$ se obtiene la ecuación del recuadro.

\begin{figure}[H]
\centering
\begin{tikzpicture}
\begin{axis}[
  width=0.78\textwidth,height=7.2cm,
  xmin=-1.1,xmax=0.75,ymin=-4.3,ymax=2.2,
  axis lines=middle,
  xlabel={$V_A$ (V)},ylabel={$I$ (mA)},
  xlabel style={at={(ticklabel* cs:1)},anchor=north west},
  ylabel style={at={(ticklabel* cs:1)},anchor=south},
  xtick={-1,-0.5,0.5},ytick={-4,-3,-2,-1,1,2},
  every axis plot/.append style={line width=1.1pt,samples=200},
  clip=true,
]
\fill[naranjaclaro,opacity=0.7] (axis cs:0,0) rectangle (axis cs:0.75,-4.3);
\fill[azulclaro,opacity=0.7] (axis cs:-1.1,0) rectangle (axis cs:0,-4.3);
\fill[verdeclaro,opacity=0.7] (axis cs:0,0) rectangle (axis cs:0.75,2.2);
\addplot[black,domain=-1.1:0.66] {1e-10*(exp(x/0.02585)-1)};
\addplot[azul,domain=-1.1:0.68] {1e-10*(exp(x/0.02585)-1)-1};
\addplot[violeta,domain=-1.1:0.69] {1e-10*(exp(x/0.02585)-1)-2};
\addplot[rojo,domain=-1.1:0.7] {1e-10*(exp(x/0.02585)-1)-3};
\node[font=\scriptsize,anchor=south west] at (axis cs:-0.85,0.03) {$G_L=0$ (oscuridad)};
\node[font=\scriptsize,azul,anchor=south west] at (axis cs:-0.85,-0.97) {$G_L=G_0$};
\node[font=\scriptsize,violeta,anchor=south west] at (axis cs:-0.85,-1.97) {$2G_0$};
\node[font=\scriptsize,rojo,anchor=south west] at (axis cs:-0.85,-2.97) {$3G_0$};
\node[font=\small,align=center,azul!80!black] at (axis cs:-0.55,-3.75) {III: \textbf{fotodetector}\\($V_A<0$, $I\propto$ luz)};
\node[font=\small,align=center,naranja!80!black] at (axis cs:0.37,-3.75) {IV: \textbf{fotovoltaico}\\(entrega potencia)};
\node[font=\small,align=center,verde!60!black] at (axis cs:0.27,1.6) {I: diodo directo\\(LED si es directo)};
\end{axis}
\end{tikzpicture}
\caption{Característica I--V de un fotodiodo para distintos niveles de iluminación
(Pierret, fig.~9.3). En el tercer cuadrante la corriente casi no depende de la tensión y es
proporcional a la luz: el diodo es una \emph{fuente de corriente controlada por luz}.
En el cuarto cuadrante $V_A>0$ e $I<0$: la potencia $P=V_AI<0$ significa que el
dispositivo \textbf{entrega} energía (celda solar).}
\label{fig:iv-pn}
\end{figure}

\subsection{Modos de operación}
\label{sec:modos}

\begin{description}[style=nextline,leftmargin=1.2em]
  \item[Modo fotoconductivo (polarización inversa, 3.\textsuperscript{er} cuadrante)]
  Es el modo normal de un fotodetector. Ventajas: la ZCE es más ancha (colecta más por
  deriva y es más rápida), la capacidad de juntura $C_j=\epsilon A/W$ es menor (menor
  constante $RC$) y la respuesta es lineal: $I\approx-I_0-I_L\approx-I_L$. Desventaja:
  circula siempre la corriente de oscuridad $I_0$ (y su ruido).
  \item[Modo fotovoltaico ($V_A>0$ fijada por la carga, 4.\textsuperscript{o} cuadrante)]
  El diodo funciona como generador: es la \textbf{celda solar} (sección~\ref{sec:solar}).
  Como fotodetector sin polarización (cortocircuito, $V_A=0$) no hay corriente de oscuridad
  y el ruido es mínimo, pero es más lento.
  \item[Polarización directa (1.\textsuperscript{er} cuadrante)] La juntura inyecta
  portadores que se recombinan. Si el material es de gap directo, emite luz: es el
  \textbf{LED} (sección~\ref{sec:led}).
\end{description}

\begin{ojo}[\diapo{10}]
\begin{itemize}
  \item ``Polarizado inverso como fotodetector (diodo PIN)'': el PIN es una
  \emph{estructura} (sección~\ref{sec:pin}), no un modo de operación. Tanto el PN como el
  PIN se usan en polarización inversa.
  \item ``Polarizado directo como LED'': un fotodiodo de Si en directa \emph{no} emite luz
  apreciable (gap indirecto). El LED es la misma física de juntura, pero hecho con un
  material de gap directo.
  \item ``La zona n, menos contaminada, tendrá una mayor \emph{longitud} de deplexión'':
  se refiere al \emph{ancho} de la zona de deplexión. Como $x_nN_D=x_pN_A$, la ZCE se
  extiende sobre todo hacia el lado menos dopado; por eso en un diodo p$^+$n casi toda la
  ZCE (y la colección por deriva) está en el lado n.
\end{itemize}
\end{ojo}

\begin{figure}[H]
\centering
\begin{minipage}{0.5\textwidth}
\centering
\begin{tikzpicture}[scale=0.9]
  % corte transversal fotodiodo p+n
  \fill[azulclaro] (0,0) rectangle (6,2.2);
  \fill[azul!35] (0,0) rectangle (6,0.5);
  \fill[rojoclaro] (1.2,2.2) -- (1.2,1.75) arc[start angle=180,end angle=270,radius=0.25] -- (4.55,1.5) arc[start angle=270,end angle=360,radius=0.25] -- (4.8,2.2) -- cycle;
  \draw (1.2,2.2) -- (1.2,1.75) arc[start angle=180,end angle=270,radius=0.25] -- (4.55,1.5) arc[start angle=270,end angle=360,radius=0.25] -- (4.8,2.2);
  \draw (0,0) rectangle (6,2.2);
  \draw (0,0.5) -- (6,0.5);
  \node at (3,1.85) {p$^+$};
  \node at (3,1.0) {n};
  \node at (3,0.25) {n$^+$};
  % SiO2 y AR
  \fill[gris!30] (0,2.2) rectangle (1.4,2.4) (4.6,2.2) rectangle (6,2.4);
  \fill[azul!50] (1.4,2.2) rectangle (4.6,2.32);
  \fill[black!70] (1.0,2.4) rectangle (1.8,2.6) (4.2,2.4) rectangle (5.0,2.6);
  \fill[black!70] (0,-0.2) rectangle (6,0);
  \draw[foton] (3,3.4) -- (3,2.4);
  \node[font=\scriptsize,anchor=west] at (6.05,2.3) {SiO$_2$};
  \draw[gris] (4.0,2.27) -- (4.6,3.1) node[above,font=\scriptsize,black] {antirreflectante};
  \draw[gris] (1.4,2.5) -- (0.6,3.1) node[above,font=\scriptsize,black] {contacto};
  \node[font=\scriptsize,anchor=west] at (6.05,-0.1) {contacto};
\end{tikzpicture}
\end{minipage}%
\begin{minipage}{0.48\textwidth}
\centering
\begin{circuitikz}[scale=0.9,transform shape]
  \draw (0,0) to[battery1,l=$V_R$] (0,3) -- (2.6,3);
  \draw (2.6,0) to[photodiode] (2.6,3);
  \draw (2.6,0) to[R,l=$R_L$] (0,0);
  \draw[flecha,azul] (3.3,2.3) -- node[right,font=\small] {$I_L$} (3.3,0.7);
\end{circuitikz}
\end{minipage}
\caption{Izquierda: corte de un fotodiodo planar p$^+$--n \diapo{10}: capa p$^+$ muy
delgada, ventana con antirreflectante y contacto en forma de anillo. Derecha: símbolo y
polarización inversa.}
\label{fig:pn-corte}
\end{figure}

\begin{ojo}[``el material P debe ser transparente'' \diapo{12}]
En un fotodiodo de Si el lado p también es Si: \emph{no puede} ser transparente a la luz
que el Si absorbe. Lo que se hace es una capa p$^+$ \textbf{muy delgada} (mucho menor que
$1/\alpha$), para que la mayor parte de la luz atraviese esa capa y se absorba en la ZCE o
cerca de ella. Una capa superior realmente transparente requiere otro material de gap
mayor (una \emph{heteroestructura}, como la ventana de InP del fotodiodo de InGaAs,
sección~\ref{sec:materiales-pin}). Tampoco es cierto que ``cada fotón al llegar a la ZCE
produce un par'' \diapo{11}: la absorción es probabilística (ley exponencial) y ocurre en
todo el material; la ZCE sólo determina \emph{cómo} se colectan los portadores.
\end{ojo}

\subsection{Respuesta espectral}

La respuesta espectral indica cómo varía $I_L$ con $\lambda$ a potencia óptica constante.

\begin{figure}[H]
\centering
\begin{tikzpicture}
\begin{axis}[
  width=0.62\textwidth,height=5.2cm,
  xmin=0.4,xmax=1.25,ymin=0,ymax=1.1,
  xlabel={$\lambda$ (\si{\micro m})},ylabel={$I_L/I_{L,\max}$},
  grid=major,grid style={gray!25},
]
\addplot[azul,line width=1.2pt,smooth] coordinates {(0.45,0.10)(0.5,0.18)(0.55,0.28)(0.6,0.40)(0.65,0.53)(0.7,0.67)(0.75,0.80)
  (0.8,0.91)(0.85,0.98)(0.88,1.0)(0.92,0.97)(0.96,0.85)(1.0,0.62)(1.04,0.35)(1.08,0.12)(1.11,0.02)(1.14,0)};
\draw[gris,dashed] (axis cs:1.107,0) -- (axis cs:1.107,1.1) node[below left,font=\scriptsize] {$\lG(\text{Si})$};
\node[font=\scriptsize,align=center] at (axis cs:0.55,0.75) {$\lambda$ corta:\\absorción en la\\superficie};
\node[font=\scriptsize,align=center] at (axis cs:1.14,0.6) {$\Eph<\EG$:\\transparente};
\end{axis}
\end{tikzpicture}
\caption{Respuesta espectral típica de un fotodiodo PN de Si (Pierret, fig.~9.4 y \diapo{9}).}
\label{fig:resp-espectral}
\end{figure}

\begin{itemize}
  \item \textbf{Corte en $\lambda$ larga:} para $\lambda>\lG$ ($\Eph<\EG$) no hay absorción.
  Para Si, $\lG=1.24/1.12\approx\SI{1.1}{\micro m}$.
  \item \textbf{Caída en $\lambda$ corta}, por dos motivos:
  \begin{enumerate}
    \item A potencia constante, cada fotón lleva más energía, así que llegan \emph{menos
    fotones por segundo}: $\Phi = P/h\nu = P\lambda/hc$.
    \item $\alpha$ crece mucho y la luz se absorbe muy cerca de la superficie, lejos de la
    ZCE y donde la recombinación superficial es alta: muchos pares se pierden.
  \end{enumerate}
\end{itemize}

\begin{ojo}[\diapo{9}]
Dice que la respuesta disminuye a $\lambda$ bajas ``debido a que se incrementa la $E_G$ al
bajar $\lG$''. Es incorrecto: $\EG$ es una propiedad fija del material. Lo que aumenta
al bajar $\lambda$ es la \textbf{energía de cada fotón}, $\Eph=1.24/\lambda$; a potencia
constante eso implica menos fotones por segundo y, por lo tanto, menos pares.
\end{ojo}

\subsection{Velocidad de respuesta}
\label{sec:pn-velocidad}

Los portadores generados en la ZCE cruzan por deriva, normalmente a la velocidad de
saturación ($\vsat\approx\SI{e7}{cm/s}$ en Si), y tardan
$t_{tr}=W/\vsat\approx\SI{10}{ps}$ para $W=\SI{1}{\micro m}$. Los generados en las zonas
neutras tienen que \emph{difundir}, y un camino de longitud $L$ por difusión tarda
\[
  t_{dif}\sim\frac{L^2}{2D}\qquad
  \text{p.\,ej. } \frac{(\SI{10}{\micro m})^2}{2\cdot\SI{35}{cm^2/s}}\approx\SI{14}{ns}.
\]
Esa ``cola'' lenta limita al fotodiodo PN a frecuencias del orden de
\textbf{decenas de MHz}. Como en un PN la ZCE es delgada ($W\lesssim\SI{1}{\micro m}$) y
$1/\alpha$ suele ser de varios micrones, la mayoría de los pares se genera en zonas
neutras. El PIN resuelve exactamente este problema.

\begin{ojo}[velocidad de saturación \diapo{10}]
La diapositiva da $\vsat\approx\SI{e5}{cm/s}$. El valor correcto para electrones y huecos
en Si es del orden de $\SI{e7}{cm/s}$ ($=\SI{e5}{m/s}$; probablemente se mezclaron las
unidades).
\end{ojo}

% =====================================================================
\section{Fotodiodo PIN}
\label{sec:pin}
% =====================================================================

\subsection{Idea: agrandar la zona de campo}

En el fotodiodo PN la ZCE es angosta y la mayoría de los pares se genera en zonas neutras
(colección lenta, por difusión, y con pérdidas por recombinación). El \textbf{PIN} intercala
entre un p$^+$ y un n$^+$ una capa \textbf{intrínseca} (en la práctica, muy poco dopada,
p$^-$ o n$^-$) de ancho $W_i$:

\begin{itemize}
  \item Como la capa i casi no tiene carga fija, se vacía por completo con muy poca tensión
  inversa (o incluso sin polarizar) y \textbf{la ZCE ocupa toda la capa i}: $W\approx W_i$,
  independiente de la tensión aplicada.
  \item Las capas p$^+$ y n$^+$, muy dopadas, casi no se vacían (la carga se concentra en
  láminas delgadas en sus bordes) y prácticamente toda la tensión cae sobre la capa i.
  \item El campo en la capa i es \textbf{aproximadamente constante}: $\mathcal E\approx
  (V_{bi}+V_R)/W_i$.
\end{itemize}

\begin{figure}[H]
\centering
\begin{tikzpicture}[x=1cm,y=1cm]
  % dispositivo
  \fill[rojoclaro] (0,0) rectangle (1.2,1);
  \fill[gris!12] (1.2,0) rectangle (7.2,1);
  \fill[azulclaro] (7.2,0) rectangle (8.6,1);
  \draw (0,0) rectangle (8.6,1);
  \draw (1.2,0) -- (1.2,1) (7.2,0) -- (7.2,1);
  \node at (0.6,0.5) {p$^+$}; \node at (4.2,0.5) {i (n$^-$)}; \node at (7.9,0.5) {n$^+$};
  \foreach \x in {0.3,0.7}{\draw[foton] (\x-0.6,1.9) -- (\x,1.05);}
  \node[colFoton,font=\small] at (-0.6,1.6) {luz};
  % rho(x)
  \begin{scope}[yshift=-2.3cm]
    \draw[flecha] (-0.3,0) -- (9,0) node[right] {$x$};
    \draw[flecha] (-0.3,-1) -- (-0.3,1.1) node[above] {$\rho(x)$};
    \fill[rojo!50] (0.95,0) rectangle (1.2,-0.9);
    \fill[azul!50] (7.2,0) rectangle (7.45,0.9);
    \fill[azul!20] (1.2,0) rectangle (7.2,0.06);
    \node[font=\scriptsize,anchor=west] at (1.25,-0.7) {$-qN_A$ (lámina delgada)};
    \node[font=\scriptsize,anchor=east] at (7.15,0.7) {$+qN_D$};
    \node[font=\scriptsize] at (4.2,0.3) {carga casi nula};
  \end{scope}
  % E(x)
  \begin{scope}[yshift=-6.0cm]
    \draw[flecha] (-0.3,0) -- (9,0) node[right] {$x$};
    \draw[flecha] (-0.3,-0.2) -- (-0.3,1.7) node[above] {$|\mathcal{E}(x)|$};
    \fill[verde!25] (0.95,0) -- (1.2,1.2) -- (7.2,1.3) -- (7.45,0) -- cycle;
    \draw[verde!60!black,line width=1.1pt] (0.6,0) -- (0.95,0) -- (1.2,1.2) -- (7.2,1.3) -- (7.45,0) -- (8.6,0);
    \draw[gris,dashed,line width=0.9pt] (3.9,0) -- (4.4,1.9) -- (4.9,0);
    \node[gris,font=\scriptsize,align=left,anchor=west] at (4.6,1.95) {PN común (sin capa i): triangular y angosto};
    \node[verde!60!black,font=\scriptsize] at (2.5,0.6) {PIN: casi constante};
    \draw[<->] (1.2,-0.35) -- node[below,font=\small] {$W\approx W_i$} (7.2,-0.35);
  \end{scope}
\end{tikzpicture}
\caption{Estructura del PIN, densidad de carga y campo eléctrico \diapo{13--14}. El área
bajo $|\mathcal E|$ es la tensión total $V_{bi}+V_R$: para la misma tensión, el PIN
reparte el campo en una zona mucho más ancha.}
\label{fig:pin-campo}
\end{figure}

\begin{figure}[H]
\centering
\begin{tikzpicture}[scale=1]
  % bandas con polarizacion inversa
  \draw[bandac] (0,3.0) -- (1.2,3.0) -- (7.2,0.6) -- (8.6,0.6) node[right,font=\small] {$\Ec$};
  \draw[bandav] (0,1.9) -- (1.2,1.9) -- (7.2,-0.5) -- (8.6,-0.5) node[right,font=\small] {$\Ev$};
  \draw[fermi] (0,1.75) node[left,font=\scriptsize,black] {$E_{Fp}$} -- (1.2,1.75);
  \draw[fermi] (7.2,0.45) -- (8.6,0.45) node[pos=0.5,below,font=\scriptsize,black] {$E_{Fn}$};
  \draw[<->,gris] (8.9,1.75) -- node[right,font=\small,black] {$q(V_{bi}+V_R)$} (8.9,0.45);
  \draw[gris,dotted] (1.2,1.75) -- (8.9,1.75) (8.6,0.45) -- (8.9,0.45);
  % generacion en i
  \draw[foton] (2.2,4.0) -- (3.4,2.65);
  \node[colFoton,font=\small] at (2.0,4.2) {$h\nu$};
  \node[elec] at (3.5,2.08) {};
  \node[hueco] at (3.5,0.98) {};
  \draw[flecha] (3.5,1.08) -- (3.5,1.98);
  \draw[flecha,colElec,line width=0.9pt] (3.65,2.0) -- (5.6,1.25) node[midway,above right,font=\scriptsize] {e$^-$ hacia n};
  \draw[flecha,colHueco,line width=0.9pt] (3.35,0.98) -- (1.9,1.55);
  \node[colHueco,font=\scriptsize] at (2.5,0.85) {hueco hacia p};
  \draw[dashed] (1.2,-0.9) -- (1.2,3.5) (7.2,-0.9) -- (7.2,3.5);
  \node[font=\small] at (0.6,-0.8) {p$^+$}; \node[font=\small] at (4.2,-0.8) {capa i};
  \node[font=\small] at (7.9,-0.8) {n$^+$};
\end{tikzpicture}
\caption{Diagrama de bandas de un PIN en inversa (Pierret, fig.~9.5b). El fotón genera un
par en la capa i; el campo (pendiente de las bandas) acelera al electrón ``cuesta abajo''
hacia n$^+$ y al hueco ``cuesta arriba'' hacia p$^+$.}
\label{fig:pin-bandas}
\end{figure}

\subsection{Ventajas: eficiencia y velocidad}

\begin{enumerate}
  \item \textbf{Ancho a medida.} Se elige $W_i\gtrsim1/\alpha$ para la $\lambda$ de trabajo,
  de modo que casi toda la luz se absorba \emph{dentro} de la zona de campo.
  \item \textbf{Velocidad.} Los portadores se colectan por deriva a $\vsat$. El límite es el
  tiempo de tránsito por la capa i:
  \begin{clave}[Frecuencia máxima del PIN]
  \[
    f_{\max}\approx\frac{1}{t_{tr}} = \frac{\vsat}{W_i}
    \qquad\text{p.\,ej. } W_i=\SI{5}{\micro m},\ \vsat=\SI{e7}{cm/s}
    \;\Rightarrow\; f_{\max}\approx\SI{20}{GHz}
  \]
  \end{clave}
  \item \textbf{Capacidad baja y constante:} $C_j=\epsilon A/W_i$ no depende de la tensión
  y es chica, lo que reduce la constante de tiempo $RC$.
\end{enumerate}

\begin{intuicion}[el compromiso del ancho $W_i$]
Si $W_i$ es grande se absorbe más luz (mejor $\eta$) pero el tránsito es más lento.
Si $W_i$ es muy chico, el tránsito es rápido pero se absorbe poca luz \emph{y} crece la
capacidad $C_j\propto1/W_i$, de modo que termina limitando $RC$. Hay un óptimo para cada
$\lambda$ y cada aplicación.
\end{intuicion}

\subsection{Eficiencia cuántica}

\begin{definicion}[eficiencia cuántica de un fotodiodo \enhoja]
\[
  \eta = \frac{\text{pares e--h colectados (que generan corriente) por segundo}}
              {\text{fotones incidentes por segundo}}
       = \frac{I_L/q}{\Popt/h\nu}
\]
En un fotodiodo ideal $\eta=1$: cada fotón produce un electrón de corriente.
\end{definicion}

Calculemos $\eta$ para un PIN \diapo{18}. Llega una potencia $P$; en la superficie se
refleja la fracción $R$, entra $(1-R)P$ y se atenúa como $e^{-\alpha x}$. La potencia
absorbida dentro de la capa de ancho $W$ es $(1-R)P\,(1-e^{-\alpha W})$, así que

\begin{clave}[Eficiencia cuántica del PIN]
\[
  \eta = \eta_i\,(1-R)\left(1-e^{-\alpha W}\right)
  \qquad\Bigl[\text{con capa p}^+\text{ de espesor } d:\
  \eta = \eta_i\,(1-R)\,e^{-\alpha d}\left(1-e^{-\alpha W}\right)\Bigr]
\]
\begin{itemize}[nosep]
  \item $(1-R)$: pérdida por reflexión $\Rightarrow$ capa antirreflectante ($R\to0$).
  \item $(1-e^{-\alpha W})$: luz que atraviesa sin absorberse $\Rightarrow$ $\alpha W\gg1$.
  \item $e^{-\alpha d}$: luz absorbida en la capa superior (se pierde) $\Rightarrow$ capa
  p$^+$ delgada o transparente.
  \item $\eta_i$ (eficiencia interna): fracción de los pares generados en la ZCE que
  efectivamente llega a los contactos (algunos se recombinan).
\end{itemize}
\end{clave}

\begin{figure}[H]
\centering
\begin{tikzpicture}
\begin{axis}[
  width=0.8\textwidth,height=5.2cm,
  xmin=-1.5,xmax=9,ymin=0,ymax=1.15,
  axis lines=left,
  xlabel={$x$},ylabel={$P(x)/P$},
  xtick={0,6},xticklabels={$0$,$W$},
  ytick={0.7,1},yticklabels={$(1-R)$,$1$},
  clip=false,
]
\fill[naranjaclaro] (axis cs:0,0) rectangle (axis cs:6,1.15);
\node[naranja!80!black,font=\small] at (axis cs:3,1.05) {zona de absorción útil (ZCE / capa i)};
\addplot[black,line width=1.2pt] coordinates {(-1.5,1)(0,1)};
\draw[gris,dashed] (axis cs:0,1) -- (axis cs:0,0.7);
\addplot[naranja,line width=1.2pt,domain=0:9,samples=80] {0.7*exp(-x/3)};
\draw[<->,azul] (axis cs:6.3,0) -- (axis cs:6.3,{0.7*exp(-2)});
\draw[azul] (axis cs:6.3,0.05) -- (axis cs:7.0,0.35) node[above,font=\small,align=center] {$(1-R)P\,e^{-\alpha W}$\\(se pierde)};
\draw[<->,verde!60!black] (axis cs:-0.5,0.7) -- node[right,font=\small] {$RP$} (axis cs:-0.5,1);
\node[font=\small,align=center] at (axis cs:3.2,0.6) {absorbida:\\$(1-R)P\,(1-e^{-\alpha W})$};
\end{axis}
\end{tikzpicture}
\caption{Potencia óptica dentro del fotodiodo. Sólo la parte absorbida en $0<x<W$ se
colecta eficientemente.}
\label{fig:pin-potencia}
\end{figure}

\subsection{Responsividad}

En la práctica interesa cuánta corriente se obtiene por cada watt de luz.

\begin{clave}[Responsividad \enhoja]
\[
  \Resp = \frac{I_L}{\Popt} = \frac{\eta\,q}{h\nu}
  = \eta\,\frac{q\lambda}{hc}
  \quad\Longrightarrow\quad
  \Resp\,[\si{A/W}] = \eta\,\frac{\lambda\,[\si{\micro m}]}{1.24}
\]
Para una señal modulada, $i(t)=\Resp\,P(t)$: el fotodiodo es un \emph{detector
cuadrático} (la corriente es proporcional a la potencia, es decir al cuadrado de la amplitud
del campo eléctrico de la onda).
\end{clave}

\begin{figure}[H]
\centering
\begin{tikzpicture}
\begin{axis}[
  width=0.6\textwidth,height=6cm,
  xmin=0.4,xmax=1.9,ymin=0,ymax=1.4,
  xlabel={$\lambda$ (\si{\micro m})},ylabel={$\Resp$ (\si{A/W})},
  grid=major,grid style={gray!25},
  legend pos=outer north east,legend style={font=\small},
  every axis plot/.append style={line width=1.1pt},
]
\addplot[black,dashed,domain=0.4:1.75] {x/1.24};
\addlegendentry{$\eta=1$ (ideal)}
\addplot[azul,smooth] coordinates {(0.4,0.12)(0.5,0.25)(0.6,0.36)(0.7,0.46)(0.8,0.55)(0.9,0.6)(0.95,0.56)(1.0,0.42)(1.05,0.22)(1.1,0.05)(1.12,0)};
\addlegendentry{Si}
\addplot[naranja,smooth] coordinates {(0.8,0.25)(1.0,0.4)(1.2,0.55)(1.4,0.66)(1.55,0.72)(1.6,0.68)(1.65,0.5)(1.7,0.25)(1.75,0.08)(1.8,0)};
\addlegendentry{Ge}
\addplot[verde,smooth] coordinates {(0.93,0)(0.97,0.3)(1.0,0.55)(1.1,0.7)(1.3,0.85)(1.55,0.95)(1.6,0.96)(1.64,0.75)(1.67,0.35)(1.70,0.05)(1.71,0)};
\addlegendentry{InGaAs (con ventana InP)}
\end{axis}
\end{tikzpicture}
\caption{Responsividad de fotodiodos típicos (valores aproximados, \diapo{21}). La recta
$\eta=1$ crece con $\lambda$ porque, a igual potencia, a mayor $\lambda$ llegan más fotones
(cada uno con menos energía). Cada curva real cae a cero en su $\lG$.}
\label{fig:responsividad}
\end{figure}

\paragraph{Por qué cae la responsividad en los extremos \diapo{21}:}
\begin{itemize}
  \item \textbf{$\lambda$ cortas} (fotones de mucha energía): $\alpha$ es tan grande que la luz
  se absorbe en la capa superior (p$^+$) antes de llegar a la ZCE; si esa capa es de otro
  material de mayor gap, para estas energías deja de ser transparente. En ambos casos
  $\eta$ baja.
  \item \textbf{$\lambda$ largas} (cerca de $\lG$): $\alpha$ es tan chico que la luz atraviesa
  toda la capa i sin absorberse ($\alpha W\ll1$), y para $\lambda>\lG$ no se absorbe en
  absoluto.
\end{itemize}

\begin{ojo}[definición de responsividad \diapo{20}]
\begin{itemize}
  \item La diapo la define como ``la corriente obtenida por unidad de \emph{energía}
  luminosa incidente''. Es por unidad de \textbf{potencia} óptica: $[\Resp]=\si{A/W}$
  (amperes por watt, no por joule).
  \item ``Tiene la misma información que $\eta$ pero además tiene información de la
  frecuencia'': se refiere a la frecuencia \emph{de la luz} $\nu=c/\lambda$ (aparece en
  $h\nu$), no a la frecuencia de modulación de la señal. La velocidad de respuesta del
  detector es otro tema (tiempo de tránsito, $RC$).
\end{itemize}
\end{ojo}

\begin{ejemplo}[diseño de un PIN de Si para \SI{0.8}{\micro m}]
A $\lambda=\SI{0.8}{\micro m}$ el Si tiene $\alpha\approx\SI{850}{cm^{-1}}$
($1/\alpha\approx\SI{12}{\micro m}$). Sea $W=\SI{10}{\micro m}$ y $\eta_i=1$.

\textbf{Sin antirreflectante} ($R=0.31$): $\alpha W=0.85$,
$\eta=0.69\,(1-e^{-0.85})=0.69\cdot0.57\approx0.39$.

\textbf{Con antirreflectante} ($R\approx0$): $\eta\approx0.57$, y
$\Resp=0.57\cdot0.8/1.24\approx\SI{0.37}{A/W}$. Con $\Popt=\SI{10}{\micro W}$:
$I_L\approx\SI{3.7}{\micro A}$.

\textbf{¿Qué $W$ hace falta para $\eta=0.9$?} $1-e^{-\alpha W}=0.9\Rightarrow
W=\ln10/\alpha\approx\SI{27}{\micro m}$. Entonces
$t_{tr}=W/\vsat=\SI{27e-4}{cm}/\SI{e7}{cm/s}=\SI{0.27}{ns}$ y
$f_{\max}\approx\SI{3.7}{GHz}$, contra $\approx\SI{10}{GHz}$ con $W=\SI{10}{\micro m}$:
el compromiso eficiencia--velocidad en números.
\end{ejemplo}

\subsection{Corriente de oscuridad y ruido}

Aunque no haya luz circula la \textbf{corriente de oscuridad} $I_{dark}$: la corriente
inversa del diodo. En un PIN de Si a temperatura ambiente suele dominar la generación
térmica en la ZCE, $I_{gen}\approx qAn_iW/(2\tau_0)$ (por eso crece con $W$), y como
depende de $n_i$ crece rápidamente con la temperatura (en Si, aproximadamente se duplica cada
\SIrange{8}{10}{\celsius}).

\begin{ojo}[``la corriente de oscuridad es continua y se puede filtrar fácilmente'' \diapo{19}]
Su \emph{valor medio} es continuo y se puede quitar con un capacitor de acople, sí. Pero
como toda corriente formada por portadores discretos, tiene \textbf{ruido de granalla}
(\emph{shot noise}) de valor eficaz $i_n=\sqrt{2q\,I_{dark}\,B}$ en un ancho de banda $B$,
que cae en la misma banda que la señal y \textbf{no se puede filtrar}. Ese ruido fija la
mínima potencia óptica detectable: por eso se busca $I_{dark}$ chica (materiales de gap
grande, buena pasivación superficial, baja temperatura).
\end{ojo}

\subsection{Materiales para PIN: Si, Ge e InGaAs}
\label{sec:materiales-pin}

Las comunicaciones por fibra óptica usan \SI{1.31}{\micro m} (mínima dispersión de la
fibra) y \SI{1.55}{\micro m} (mínima atenuación). Ambas están por encima de
$\lG(\text{Si})=\SI{1.1}{\micro m}$, así que hace falta otro material.

El material elegido es la aleación \textbf{In$_{0.53}$Ga$_{0.47}$As}: gap directo de
\SI{0.75}{\eV} ($\lG=\SI{1.65}{\micro m}$) y, con esa composición, la \emph{misma constante
de red} que el InP (\diapo{15}, Pierret fig.~9.6). Que las constantes de red coincidan
(\emph{lattice matching}) permite crecer capas cristalinas de buena calidad, sin
dislocaciones, sobre sustratos de InP. En el gráfico $\EG$ vs.\ constante de red, cada
aleación ternaria ocupa la curva que une a sus dos compuestos binarios; las aleaciones
que sirven sobre un sustrato dado son las que caen sobre la vertical de su constante de red.

\begin{figure}[H]
\centering
\begin{tikzpicture}[scale=0.95]
  \fill[black!70] (0,4.6) rectangle (1.3,4.85) (3.7,4.6) rectangle (5,4.85);
  \fill[gris!40] (0,4.35) rectangle (5,4.6);
  \fill[rojoclaro] (1.1,4.35) -- (1.1,4.1) -- (3.9,4.1) -- (3.9,4.35) -- cycle;
  \fill[violetaclaro] (0,2.7) rectangle (5,4.35);
  \fill[azulclaro] (0,1.9) rectangle (5,2.7);
  \fill[azul!30] (0,0.4) rectangle (5,1.9);
  \fill[black!70] (0,0.15) rectangle (1.8,0.4) (3.2,0.15) rectangle (5,0.4);
  \draw (0,0.4) rectangle (5,4.6) (0,1.9) -- (5,1.9) (0,2.7) -- (5,2.7) (0,4.35) -- (5,4.35);
  \draw[densely dashed] (1.1,4.35) -- (1.1,4.1) -- (3.9,4.1) -- (3.9,4.35);
  \node[font=\small] at (2.5,4.22) {p$^+$};
  \node[font=\small] at (2.5,3.4) {i (n$^-$) -- In$_{0.53}$Ga$_{0.47}$As};
  \node[font=\small] at (2.5,2.3) {n -- InP (buffer)};
  \node[font=\small] at (2.5,1.15) {n$^+$ -- InP (sustrato)};
  \draw[foton] (2.5,-0.8) -- (2.5,0.35);
  \node[colFoton,font=\small,anchor=west] at (2.6,-0.6) {$h\nu$ (1.3 o \SI{1.55}{\micro m})};
  \draw[gris] (5,4.5) -- (5.6,4.9) node[right,font=\scriptsize,black] {nitruro de silicio (pasivación)};
  \draw[gris] (5,4.4) -- (5.6,4.4) node[right,font=\scriptsize,black] {InGaAsP (capa de cobertura)};
  \node[font=\scriptsize,align=left,anchor=west] at (5.3,1.15) {InP: $\lG=\SI{0.92}{\micro m}$\\$\Rightarrow$ \textbf{transparente} a\\1.3 y \SI{1.55}{\micro m} (``ventana'')};
\end{tikzpicture}
\caption{Fotodiodo PIN de InGaAs (Pierret, fig.~9.7 y \diapo{16}). La luz entra
\emph{por el sustrato}: el InP, de gap mayor, es transparente a esas longitudes de onda y
la luz llega intacta a la capa i de InGaAs, que la absorbe. Este sí es un caso en el que
la capa que atraviesa la luz es realmente transparente.}
\label{fig:ingaas}
\end{figure}

\subsection{Fotodiodo de avalancha (APD)}
\label{sec:apd}

Complemento (Pierret 9.2.2): un APD es un fotodiodo (PN o PIN) polarizado \emph{cerca de la
tensión de ruptura}. Los portadores fotogenerados ganan tanta energía en el campo que
producen nuevos pares por \textbf{ionización por impacto} (figura~\ref{fig:rg}f), y cada
par fotogenerado se multiplica por un factor $M$ (de 10 a 100 típicamente):
$I_L^{APD}=M\,I_L$. Llevan un \emph{anillo de guarda} para que la ruptura no ocurra antes
en los bordes de la juntura. La ventaja es que la ganancia ocurre \emph{antes} del
amplificador, así que mejora la relación señal/ruido cuando domina el ruido del circuito
receptor; la desventaja es que la multiplicación agrega su propio ruido y requiere
tensiones altas y estables.

% =====================================================================
\section{Celdas solares}
\label{sec:solar}
% =====================================================================

Una celda solar es, en esencia, un fotodiodo PN de área grande que trabaja en el
\textbf{cuarto cuadrante}: tensión positiva y corriente negativa, es decir, entrega
potencia a una carga. La física es la misma que en la sección~\ref{sec:pn}; lo que cambia
son los objetivos de diseño.

\begin{table}[H]
\centering\small
\caption{Fotodiodo vs.\ celda solar \diapo{27}.}
\label{tab:fd-vs-celda}
\begin{tabular}{p{0.28\textwidth}p{0.3\textwidth}p{0.3\textwidth}}
\toprule
 & \textbf{Fotodiodo} & \textbf{Celda solar} \\
\midrule
Luz            & haz de $\lambda$ definida, potencia chica & espectro solar completo, mucha potencia \\
Área           & chica (mínima capacidad, velocidad) & muy grande (captar potencia) \\
Polarización   & inversa (3.\textsuperscript{er} cuadrante) & ninguna; la carga fija el punto (4.\textsuperscript{o} cuadrante) \\
Figura de mérito & $\eta$, $\Resp$, velocidad, ruido & rendimiento de potencia $\eta=P_m/P_{in}$ \\
\bottomrule
\end{tabular}
\end{table}

\subsection{Característica I--V y parámetros}

Con la convención de la celda como \emph{generador} (corriente positiva saliendo por el
terminal p hacia la carga) la ecuación del fotodiodo se escribe
\[
  I = I_L - I_0\left(e^{V/V_T}-1\right),
  \qquad
  P = V\,I = I_L V - I_0 V\left(e^{V/V_T}-1\right).
\]
Es la misma ecuación de la sección~\ref{sec:pn} con el signo de $I$ cambiado: por eso la
curva se dibuja ``dada vuelta'' (Pierret, fig.~9.9).

\begin{clave}[Parámetros de la celda solar \enhoja]
\begin{itemize}[nosep]
  \item \textbf{Corriente de cortocircuito} ($V=0$): $\Isc = I_L$, proporcional a la
  irradiancia.
  \item \textbf{Tensión de circuito abierto} ($I=0$):
  \[
    \Voc = V_T\ln\!\left(1+\frac{I_L}{I_0}\right)\approx V_T\ln\frac{I_L}{I_0}.
  \]
  Crece sólo logarítmicamente con la luz y \emph{disminuye} con la temperatura (porque
  $I_0\propto n_i^2$ crece muy rápido); en Si, unos \SI{-2}{mV/\celsius}.
  \item \textbf{Punto de máxima potencia} $(V_m,I_m)$: donde $\dd P/\dd V=0$.
  \item \textbf{Factor de forma} (\emph{fill factor}):
  \[
    FF = \frac{P_m}{\Isc\Voc} = \frac{V_mI_m}{\Isc\Voc} < 1 .
  \]
  \item \textbf{Rendimiento} (eficiencia de conversión):
  \[
    \eta = \frac{P_m}{P_{in}} = \frac{FF\,\Isc\,\Voc}{P_{in}} .
  \]
\end{itemize}
\end{clave}

\begin{figure}[H]
\centering
\begin{tikzpicture}
\begin{axis}[
  width=0.78\textwidth,height=6.8cm,
  xmin=0,xmax=0.7,ymin=0,ymax=4.2,
  axis lines=left,
  xtick={0,0.1,0.2,0.3,0.4},ytick={0,1,2},
  xlabel={$V$ (V)},ylabel={$I$ (A)},
  ylabel style={azul},
  every axis plot/.append style={samples=300},
  clip=false,
]
\fill[verde!20] (axis cs:0,0) rectangle (axis cs:0.5475,3.342);
\draw[verde!60!black,pattern=north east lines,pattern color=verde!50] (axis cs:0,0) rectangle (axis cs:0.5475,3.342);
\draw[gris,dashed] (axis cs:0,3.5) -- (axis cs:0.6276,3.5) -- (axis cs:0.6276,0);
\addplot[azul,line width=1.3pt,domain=0:0.6276] {3.5-1e-10*(exp(x/0.02585)-1)};
\fill[azul] (axis cs:0.5475,3.342) circle (2pt);
\node[font=\small,anchor=south west] at (axis cs:0.5475,3.342) {$P_m$};
\node[font=\small,anchor=south east] at (axis cs:0,3.45) {$\Isc$};
\node[font=\small,anchor=north east] at (axis cs:0,3.40) {$I_m$};
\node[font=\small,anchor=north] at (axis cs:0.5475,0) {$V_m$};
\node[font=\small,anchor=north] at (axis cs:0.6276,0) {$\Voc$};
\node[verde!40!black,font=\small,align=center,fill=white,inner sep=2pt] at (axis cs:0.26,1.6) {rectángulo de\\máxima potencia\\$P_m=V_mI_m$};
\node[gris,font=\scriptsize,align=center] at (axis cs:0.3,3.85) {rectángulo $\Isc\Voc$\quad($FF$ = verde / punteado)};
\end{axis}
\begin{axis}[
  width=0.78\textwidth,height=6.8cm,
  xmin=0,xmax=0.7,ymin=0,ymax=2.2,
  axis x line=none,axis y line*=right,
  ylabel={$P$ (W)},ylabel style={rojo},
  every axis plot/.append style={samples=300},
]
\addplot[rojo,line width=1.1pt,dashed,domain=0:0.6276] {x*(3.5-1e-10*(exp(x/0.02585)-1))};
\end{axis}
\end{tikzpicture}
\caption{Curva I--V (azul, eje izquierdo) y potencia entregada $P=VI$ (rojo, eje derecho)
de la celda del ejemplo de abajo. El factor de forma mide qué tan ``cuadrada'' es la curva.}
\label{fig:iv-celda}
\end{figure}

\paragraph{Punto de máxima potencia.} Derivando $P$:
\[
  \frac{\dd P}{\dd V}=0
  \;\Rightarrow\;
  \left(1+\frac{V_m}{V_T}\right)e^{V_m/V_T} = 1+\frac{I_L}{I_0}
  \;\Rightarrow\;
  V_m = \Voc - V_T\ln\!\left(1+\frac{V_m}{V_T}\right).
\]
Es una ecuación trascendente; se resuelve iterando (empezando con $V_m=\Voc$). Una
aproximación empírica muy usada (Green) para el factor de forma ideal es
\[
  FF\approx\frac{v_{oc}-\ln(v_{oc}+0.72)}{v_{oc}+1},\qquad v_{oc}=\Voc/V_T .
\]

\begin{ejemplo}[celda de Si de \SI{10}{cm}$\times$\SI{10}{cm}]
$A=\SI{100}{cm^2}$, $J_{sc}=\SI{35}{mA/cm^2}$ bajo \SI{1000}{W/m^2}
($P_{in}=\SI{100}{mW/cm^2}\cdot\SI{100}{cm^2}=\SI{10}{W}$), $I_0=\SI{e-10}{A}$, $T=\SI{300}{K}$.
\begin{itemize}[nosep]
  \item $\Isc=I_L=\SI{3.5}{A}$.
  \item $\Voc=0.02585\ln(3.5/10^{-10})=0.02585\cdot24.28=\SI{0.628}{V}$.
  \item Iterando: $V_m\approx0.628-0.02585\ln(1+0.628/0.02585)\approx\SI{0.545}{V}$; una
  segunda vuelta da $V_m\approx\SI{0.547}{V}$, $I_m=3.5-10^{-10}e^{0.547/0.02585}\approx\SI{3.34}{A}$.
  \item $P_m\approx\SI{1.83}{W}$, \quad $FF=1.83/(3.5\cdot0.628)\approx0.83$,
  \quad $\eta=1.83/10\approx\SI{18}{\percent}$.
\end{itemize}
\end{ejemplo}

\begin{ojo}[\diapo{27}]
\begin{itemize}
  \item En $P = IV = I_sV(e^{V/V_T}-1)-I_LV$ está usada la convención del diodo
  (corriente positiva entrando por p); con ella $P<0$ en el cuarto cuadrante significa
  potencia \emph{entregada}. Es correcto, pero conviene no mezclar convenciones.
  \item ``$\eta$ está entre el 30 y 35\,\% con semiconductores III-V'': eso vale para celdas
  \emph{multijuntura} (ver más abajo; los récords superan el \SI{39}{\percent}, y bajo
  concentración el \SI{47}{\percent}). Una celda III-V de una sola juntura (GaAs) llega a
  $\approx\SI{29}{\percent}$.
  \item ``Las celdas de Si tienen menor rendimiento por ser de banda indirecta'': el gap
  del Si (\SI{1.12}{\eV}) está prácticamente en el óptimo teórico (figura~\ref{fig:sq}). Lo
  que implica el gap indirecto es un $\alpha$ chico, que obliga a usar obleas gruesas
  ($\sim\SI{150}{\micro m}$) de muy alta calidad (largas $L_n$) y técnicas de atrapamiento
  de luz. El récord de Si de una juntura ronda el \SI{27}{\percent}.
\end{itemize}
\end{ojo}

\subsection{Circuito equivalente y resistencias parásitas}

\begin{figure}[H]
\centering
\begin{minipage}{0.45\textwidth}
\centering
\begin{circuitikz}[scale=0.85,transform shape]
  \draw (0,0) to[I,l=$I_L$,invert] (0,3)
        -- (2,3) to[D*,l=$D$,-*] (2,0)
        (2,3) -- (4,3) to[R,l=$R_{sh}$] (4,0)
        (4,3) to[R,l=$R_s$,-o] (6.5,3)
        (6.5,3) to[R,l=$R_L$] (6.5,0) -- (0,0);
  \draw (2,0) -- (4,0);
  \draw[flecha] (5.4,3.4) -- node[above,font=\small] {$I$} (6.2,3.4);
\end{circuitikz}
\end{minipage}%
\begin{minipage}{0.53\textwidth}
\centering
\begin{tikzpicture}
\begin{axis}[
  width=\textwidth,height=5.5cm,
  xmin=0,xmax=0.66,ymin=0,ymax=3.9,
  xlabel={$V$ (V)},ylabel={$I$ (A)},axis lines=left,
  legend pos=south west,legend columns=2,legend style={font=\scriptsize},
  every axis plot/.append style={line width=1.1pt},
]
\addplot[black,domain=0:0.6276,samples=200] {3.5-1e-10*(exp(x/0.02585)-1)};
\addlegendentry{ideal}
\addplot[azul] coordinates {(0,3.5)(0.3,3.5)(0.36,3.498)(0.4,3.492)(0.42,3.483)(0.44,3.464)(0.46,3.424)(0.48,3.346)(0.5,3.201)(0.52,2.961)(0.54,2.61)(0.56,2.149)(0.58,1.594)(0.6,0.964)(0.62,0.275)(0.6276,0)};
\addlegendentry{$R_s=\SI{20}{m\ohm}$}
\addplot[violeta] coordinates {(0,3.5)(0.2,3.5)(0.26,3.498)(0.3,3.491)(0.32,3.48)(0.34,3.459)(0.36,3.417)(0.38,3.344)(0.4,3.229)(0.42,3.069)(0.44,2.868)(0.46,2.631)(0.48,2.368)(0.5,2.084)(0.52,1.784)(0.54,1.47)(0.56,1.147)(0.58,0.815)(0.6,0.476)(0.62,0.132)(0.6276,0)};
\addlegendentry{$R_s=\SI{50}{m\ohm}$}
\addplot[naranja] coordinates {(0,3.5)(0.1,3.3)(0.2,3.1)(0.3,2.9)(0.4,2.699)(0.44,2.618)(0.48,2.528)(0.5,2.475)(0.52,2.406)(0.54,2.302)(0.56,2.124)(0.58,1.785)(0.6,1.097)(0.617,0)};
\addlegendentry{$R_{sh}=\SI{0.5}{\ohm}$}
\end{axis}
\end{tikzpicture}
\end{minipage}
\caption{Izquierda: circuito equivalente de una celda real. Derecha: efecto de las
resistencias parásitas sobre la celda del ejemplo. $R_s$ (contactos, dedos, emisor)
inclina la curva cerca de $\Voc$; una $R_{sh}$ baja (fugas por bordes o defectos) la
inclina cerca de $\Isc$. Ambas reducen el $FF$.}
\label{fig:circuito-celda}
\end{figure}

Con resistencias parásitas la ecuación queda implícita:
\[
  I = I_L - I_0\left(e^{(V+IR_s)/V_T}-1\right) - \frac{V+IR_s}{R_{sh}} .
\]
Una celda buena tiene $R_s$ lo más chica posible y $R_{sh}$ lo más grande posible.
Notar lo sensible que es: con $R_s=\SI{50}{m\ohm}$ y $I\approx\SI{3.3}{A}$ se pierden
$\approx\SI{170}{mV}$ en la resistencia, casi un tercio de $V_m$.

\subsection{El espectro solar y los límites del rendimiento}

Fuera de la atmósfera el Sol entrega $\approx\SI{1360}{W/m^2}$ (espectro \textbf{AM0}). Al
atravesar la atmósfera se pierde parte (absorción de O$_3$, H$_2$O, CO$_2$): el espectro
estándar terrestre \textbf{AM1.5} se normaliza a \SI{1000}{W/m^2} (\SI{100}{mW/cm^2}).
\emph{AM} (\emph{air mass}) indica cuántas ``atmósferas'' atraviesa la luz; 1.5 corresponde
al Sol a unos \ang{48} del cenit.

Dos mecanismos fundamentales limitan cualquier celda de una sola juntura:
\begin{enumerate}
  \item \textbf{Pérdida por transparencia:} los fotones con $\Eph<\EG$ ($\lambda>\lG$) no se
  absorben. En Si se pierde así $\sim\SI{20}{\percent}$ de la energía incidente; en GaAs,
  $\sim\SI{35}{\percent}$.
  \item \textbf{Pérdida por termalización:} un fotón con $\Eph>\EG$ genera \emph{un} par; el
  exceso $\Eph-\EG$ se pierde como calor en $\sim$\si{ps} (los portadores ``caen'' al
  borde de banda). De cada fotón sólo se aprovecha $\EG$.
\end{enumerate}

\begin{figure}[H]
\centering
\begin{tikzpicture}
\begin{axis}[
  width=0.8\textwidth,height=6.2cm,
  xmin=0.25,xmax=2.5,ymin=0,ymax=2200,
  xlabel={$\lambda$ (\si{\micro m})},ylabel={Irradiancia espectral (\si{W/m^2.\micro m})},
  axis lines=left,
]
\addplot[name path=sol,black,line width=1.1pt,domain=0.25:2.5,samples=200] {8101/(x^5*(exp(2.49/x)-1))};

\addplot[name path=util,draw=none,domain=0.25:1.107,samples=150] {8101/(x^5*(exp(2.49/x)-1))*x/1.107};
\addplot[name path=eje,draw=none,domain=0.25:2.5] {0};
\addplot[verde!50] fill between[of=util and eje,soft clip={domain=0.25:1.107}];
\node[verde!40!black,font=\small] at (axis cs:0.62,350) {aprovechable};
\addplot[naranja!50] fill between[of=sol and util,soft clip={domain=0.25:1.107}];
\node[naranja!60!black,font=\small] at (axis cs:0.53,1250) {termalización};
\addplot[azul!30] fill between[of=sol and eje,soft clip={domain=1.107:2.5}];
\node[azul!70!black,font=\small,align=center] (na) at (axis cs:1.85,750) {no absorbida\\($\lambda>\lG$)};
\draw[azul!70!black,-{Stealth[length=4pt]}] (na) -- (axis cs:1.45,120);
\node[font=\small,align=center] (sl) at (axis cs:1.75,1650) {Sol: cuerpo negro\\a \SI{5800}{K} ($\approx$AM0)};
\draw[-{Stealth[length=4pt]}] (sl) -- (axis cs:0.64,1650);
\draw[gris,dashed] (axis cs:1.107,0) -- (axis cs:1.107,1000) node[above,font=\scriptsize,black] {$\lG(\text{Si})$};
\draw[gris,dashed] (axis cs:0.873,0) -- (axis cs:0.873,1200) node[above,font=\scriptsize,black] {$\lG(\text{GaAs})$};
\end{axis}
\end{tikzpicture}
\caption{Espectro solar aproximado (cuerpo negro; el real tiene bandas de absorción
atmosféricas, ver Pierret fig.~9.10 y \diapo{28}) y su reparto para una celda de Si: de un
fotón de longitud de onda $\lambda<\lG$ sólo se aprovecha la fracción
$\EG/\Eph=\lambda/\lG$.}
\label{fig:espectro-solar}
\end{figure}

\begin{figure}[H]
\centering
\begin{tikzpicture}
\begin{axis}[
  width=0.62\textwidth,height=5.4cm,
  xmin=0.4,xmax=3.6,ymin=0,ymax=50,
  xlabel={$\EG$ (\si{\eV})},ylabel={rendimiento máximo (\%)},
  grid=major,grid style={gray!25},axis lines=left,
]
\addplot[azul,line width=1.2pt,smooth] coordinates {(0.5,31.7)(0.7,38.9)(0.9,42.9)(1.0,43.7)(1.1,43.9)(1.2,43.5)(1.4,41.3)(1.6,37.9)(1.8,33.8)(2.0,29.4)(2.2,25.0)(2.5,19.1)(3.0,11.3)(3.5,6.3)};
\node[azul,font=\scriptsize,anchor=south west] at (axis cs:1.5,41) {sólo pérdidas espectrales};
\draw[rojo,dashed,line width=1pt] (axis cs:1.34,7) -- (axis cs:1.34,33.7);
\fill[rojo] (axis cs:1.34,33.7) circle (2pt);
\node[rojo,font=\scriptsize,align=center] (sq) at (axis cs:3.0,32) {Shockley--Queisser:\\\SI{33.7}{\percent} en \SI{1.34}{\eV}};
\draw[rojo,-{Stealth[length=4pt]}] (sq) -- (axis cs:1.40,33.7);
\draw[gris] (axis cs:0.66,0) -- (axis cs:0.66,4) node[above,font=\scriptsize,black] {Ge};
\draw[gris] (axis cs:1.12,0) -- (axis cs:1.12,4) node[above,font=\scriptsize,black] {Si};
\draw[gris] (axis cs:1.42,0) -- (axis cs:1.42,4) node[above,font=\scriptsize,black] {GaAs};
\draw[gris] (axis cs:2.26,0) -- (axis cs:2.26,4) node[above,font=\scriptsize,black] {GaP};
\end{axis}
\end{tikzpicture}
\caption{Rendimiento máximo de una celda de una juntura en función de $\EG$. La curva azul
considera sólo las dos pérdidas espectrales (calculada con el espectro de cuerpo negro): un
gap chico absorbe casi todo pero termaliza mucho; un gap grande termaliza poco pero deja
pasar la mayoría de los fotones. Agregando la recombinación radiativa inevitable se obtiene
el límite de Shockley--Queisser, $\approx\SI{33.7}{\percent}$ en
$\EG\approx\SI{1.34}{\eV}$. Si y GaAs están cerca del óptimo.}
\label{fig:sq}
\end{figure}

\subsection{Diseño de la celda de Si}

Además de las pérdidas fundamentales hay pérdidas tecnológicas que el diseño ataca
(Pierret 9.3.2, figura de la \diapo{30}):

\begin{itemize}
  \item \textbf{Sombreado vs.\ resistencia serie:} el contacto frontal se hace con
  \emph{dedos} (\emph{fingers}) metálicos finos unidos por \emph{bus bars}. Dedos más
  anchos o juntos bajan $R_s$ pero tapan más luz.
  \item \textbf{Reflexión:} una superficie de Si plana refleja $\sim\SI{30}{\percent}$. Con
  \textbf{texturizado} (pirámides, por ataque químico anisotrópico) y antirreflectante se
  baja a menos del \SI{1}{\percent}.
  \item \textbf{Atrapamiento de luz:} el contacto trasero (con un óxido) actúa como espejo;
  la luz infrarroja, poco absorbida, rebota varias veces dentro de la celda.
  \item \textbf{Recombinación:} se usan obleas con vidas medias tan largas que $L_n$ supera
  el espesor de la celda, y superficies pasivadas (óxido), para que casi todo par
  generado llegue a la juntura.
\end{itemize}

\paragraph{Respuesta espectral de la celda \diapo{29}.} Cada región aporta en un rango
distinto: los fotones de alta energía (azul) se absorben en el emisor n$^+$ superficial,
donde la recombinación (superficial y por el alto dopaje) es alta; los de menor energía
llegan a la base p, que aporta la mayor parte de la corriente. La ZCE es delgada y aporta
poco. Por eso las mejoras apuntan a la superficie (``celda violeta'' con emisor muy
delgado) y a la cara trasera (BSF).

\begin{figure}[H]
\centering
\begin{tikzpicture}[scale=0.95]
  % n+ | p | p+
  \draw[bandac] (0,2.4) -- (1,2.4) to[out=0,in=180] (2.2,3.4) -- (7,3.4) to[out=0,in=180] (7.6,3.6) -- (8.8,3.6) node[right,font=\small] {$\Ec$};
  \draw[bandav] (0,1.3) -- (1,1.3) to[out=0,in=180] (2.2,2.3) -- (7,2.3) to[out=0,in=180] (7.6,2.5) -- (8.8,2.5) node[right,font=\small] {$\Ev$};
  \draw[fermi] (0,2.45) node[left,font=\scriptsize,black] {$\EF$} -- (8.8,2.45);
  \draw[dashed,gris] (1.6,0.7) -- (1.6,4.3) (7.3,0.7) -- (7.3,4.3);
  \node[font=\small] at (0.6,0.8) {n$^+$}; \node[font=\small] at (4.5,0.8) {p (base)};
  \node[font=\small] at (8.1,0.8) {p$^+$ (BSF)};
  % electron rebotando
  \node[elec] at (6.3,3.4) {};
  \draw[flecha,colElec] (6.45,3.5) -- (7.1,3.55);
  \draw[flecha,colElec] (7.1,3.65) to[out=150,in=30] (5.8,3.65);
  \node[colElec,font=\scriptsize,align=center] at (6.5,4.25) {el escalón refleja\\a los electrones};
  \node[elec] at (3.2,3.4) {};
  \draw[flecha,colElec] (3.05,3.4) -- (2.0,2.95);
  \node[colElec,font=\scriptsize] at (2.9,3.8) {colectado};
\end{tikzpicture}
\caption{Celda con \emph{Back Surface Field} (BSF, \diapo{31}). Una capa p$^+$ en la cara
trasera crea un escalón en la banda de conducción (un pequeño campo) que rechaza a los
electrones minoritarios de la base y les impide recombinarse en el contacto trasero.
Resultado: más corriente (mejor respuesta al rojo, que se absorbe en profundidad) y mayor
$\Voc$.}
\label{fig:bsf}
\end{figure}

\subsection{Heterojunturas, ventanas y multijuntura}

Una \textbf{heterojuntura} une dos semiconductores distintos (con $\EG$ distinto,
idealmente con la misma constante de red). Los escalones $\Delta E_c$ y $\Delta E_v$ en las
bandas permiten guiar portadores y separar el espectro.

\begin{figure}[H]
\centering
\begin{tikzpicture}[scale=0.9]
  % semiconductor A (gap grande) | B (gap chico)
  \draw[bandac] (0,4.2) -- (2.6,4.2) to[out=0,in=180] (3.5,4.6) -- (3.5,3.6) to[out=0,in=180] (4.6,2.9) -- (8,2.9) node[right,font=\small] {$E_{c2}$};
  \draw[bandav] (0,0.6) -- (2.6,0.6) to[out=0,in=180] (3.5,1.0) -- (3.5,1.75) to[out=0,in=180] (4.6,1.05) -- (8,1.05) node[right,font=\small] {$E_{v2}$};
  \node[colEc,font=\small,anchor=east] at (0,4.2) {$E_{c1}$};
  \node[colEv,font=\small,anchor=east] at (0,0.6) {$E_{v1}$};
  \draw[fermi] (-0.6,1.2) node[left,font=\scriptsize,black] {$\EF$} -- (8,1.2);
  \draw[<->,gris] (1.2,0.6) -- node[right,font=\small,black] {$E_{G1}$} (1.2,4.2);
  \draw[<->,gris] (6.5,1.05) -- node[right,font=\small,black] {$E_{G2}$} (6.5,2.9);
  \draw[<->,gris] (3.75,3.6) -- node[right,font=\scriptsize,black] {$\Delta E_c$} (3.75,4.6);
  \draw[<->,gris] (3.75,1.0) -- node[right,font=\scriptsize,black] {$\Delta E_v$} (3.75,1.75);
  \node[font=\small] at (1.5,5.1) {semiconductor A (ventana)};
  \node[font=\small] at (6,5.1) {semiconductor B};
  \draw[foton] (-1.6,2.6) -- (-0.1,2.6);
  \node[colFoton,font=\scriptsize,align=center] at (-1,3.1) {luz};
\end{tikzpicture}
\caption{Heterojuntura con $E_{G1}>E_{G2}$ \diapo{32}. Los fotones con $\Eph>E_{G1}$ se
absorben en A; los que tienen $E_{G2}<\Eph<E_{G1}$ atraviesan A (que para ellos es
transparente) y se absorben en B, cerca de la juntura.}
\label{fig:hetero}
\end{figure}

\begin{itemize}
  \item \textbf{Celda con ventana (``heteroface'', \diapo{33}):} sobre una homojuntura de GaAs
  ($E_{G2}$) se crece una capa delgada de Al$_x$Ga$_{1-x}$As de gap mayor ($E_{G1}$). Esa
  ventana deja pasar casi toda la luz útil y \emph{pasiva} la superficie: los electrones
  generados en el GaAs no pueden llegar a la superficie (donde se recombinarían) porque
  el escalón $\Delta E_c$ los refleja. Mejora la respuesta al azul. Cuanto más Al ($x$),
  mayor $E_{G1}$ y más se extiende la respuesta hacia energías altas.
  \item \textbf{Multijuntura (tándem):} se apilan celdas de gap decreciente (la de mayor gap
  arriba). Cada una aprovecha la porción del espectro justo por encima de su gap, reduciendo a
  la vez la termalización y la transparencia. Es la forma de superar el límite de
  Shockley--Queisser: son las celdas III-V de \SIrange{30}{47}{\percent}, usadas en
  satélites y en sistemas de concentración.
  \item \textbf{Concentradores:} lentes o espejos concentran la luz sobre una celda chica y
  cara. Al aumentar $I_L$, aumenta $\Voc$ (logarítmicamente) y la eficiencia; el problema es
  disipar el calor.
  \item \textbf{Película delgada:} Si amorfo (a-Si), CdTe, CIGS. Son de gap directo o casi
  (alto $\alpha$), así que alcanzan pocos micrones de material: más baratas y de gran área,
  pero con menor rendimiento.
\end{itemize}

\subsection{Tecnologías comerciales de paneles}

Un panel (módulo) conecta celdas en serie para sumar tensión (una celda de Si da
$\approx\SI{0.5}{V}$ en $V_m$; un panel de 60--72 celdas da \SIrange{30}{40}{V}). Los datos
de catálogo se miden en \textbf{condiciones estándar (STC)}: \SI{1000}{W/m^2}, espectro
AM1.5 y celda a \SI{25}{\celsius}. El rendimiento del panel es
$\eta=P_{\max}/(\text{área}\cdot\SI{1000}{W/m^2})$.

\begin{table}[H]
\centering\small
\caption{Comparación de tecnologías de paneles (valores de las \diapo{37--42}, típicos de
2020--2021). ``C.\,temp.'': coeficiente de temperatura de la potencia; ``Degr.'':
pérdida de potencia en 25 años.}
\label{tab:paneles}
\begin{tabularx}{\textwidth}{lccccX}
\toprule
Tecnología & Oblea & $P_{\max}$ & $\eta$ & C.\,temp. & Característica \\
           &       &           &        & (\%/\si{\celsius}) & \\
\midrule
Policristalino           & p & \SI{285}{W} & 16--17\,\% & $-0.40$ & fabricación más barata; degr.\ 20\,\% \\
Mono PERC \emph{half cell} & p & \SI{345}{W} & 18--20\,\% & $-0.35$ & capa pasivante + reflector trasero; celdas partidas \\
Shingled mono PERC       & p & \SI{390}{W} & 19--20.5\,\% & $-0.28$ & tiras solapadas, sin cintas: menos sombra y microrroturas \\
Mono PERC half cell MBB  & p & \SI{405}{W} & 19--21\,\% & $-0.35$ & múltiples \emph{bus bars}, obleas grandes M10/M12 \\
HJT (heterojuntura)      & n & \SI{380}{W} & 20--22\,\% & $-0.26$ & c-Si + capas de a-Si; ideal para climas muy cálidos; degr.\ 8\,\% \\
IBC                      & n & \SI{400}{W} & 20--23\,\% & $-0.27$ & todos los contactos atrás: cara frontal sin sombra; degr.\ 8\,\% \\
\bottomrule
\end{tabularx}
\end{table}

\begin{itemize}
  \item \textbf{PERC} (\emph{Passivated Emitter and Rear Cell}): agrega una capa dieléctrica
  en la cara trasera que pasiva la superficie (menos recombinación) y refleja la luz que no
  se absorbió para darle una segunda pasada. Es la versión industrial de las ideas de
  la figura de la \diapo{30}.
  \item \textbf{Half cell}: cada celda se corta en dos. La corriente por celda se reduce a la
  mitad, y como las pérdidas óhmicas son $I^2R$, se reducen a la cuarta parte por celda.
  Además el panel se arma como dos mitades en paralelo, que toleran mejor las sombras
  parciales.
  \item \textbf{Oblea tipo n} (HJT, IBC): no sufre la degradación inducida por la luz
  (LID) asociada al boro de las obleas p, de ahí la menor degradación a 25 años.
\end{itemize}

\begin{ojo}[\diapo{38}]
``La célula partida ayuda a reducir las pérdidas por resistencia \emph{térmica}'': las
pérdidas que se reducen son las \textbf{óhmicas} ($I^2R$ en dedos, cintas y contactos), al
reducirse la corriente a la mitad. La menor temperatura de operación es una consecuencia,
no la causa.
\end{ojo}

\begin{ejemplo}[efecto de la temperatura, \diapo{37}]
Un panel policristalino de \SI{285}{W} (STC) con coeficiente \SI{-0.40}{\percent/\celsius}
trabaja con las celdas a \SI{50}{\celsius} (normal en verano):
$\Delta T=\SI{25}{\celsius}$ $\Rightarrow$ pérdida $25\cdot\SI{0.40}{\percent}=\SI{10}{\percent}$,
es decir $\SI{28.5}{W}$: entrega $\approx\SI{256}{W}$. Un panel HJT
(\SI{-0.26}{\percent/\celsius}) perdería sólo \SI{6.5}{\percent}. La causa física es la
caída de $\Voc$ con la temperatura ($I_0\propto n_i^2$ crece exponencialmente).
\end{ejemplo}

% =====================================================================
\section{Diodos emisores de luz (LED)}
\label{sec:led}
% =====================================================================

\subsection{Principio de funcionamiento}

El LED es el proceso inverso al fotodiodo. Al polarizar en \textbf{directa} una juntura PN,
la barrera baja y se \emph{inyectan} minoritarios: electrones del lado n hacia el lado p y
huecos del lado p hacia el lado n. Esos portadores en exceso se recombinan a pocas
longitudes de difusión de la juntura. Si la recombinación es \textbf{banda a banda
radiativa}, cada recombinación emite un fotón de energía $h\nu\approx\EG$.

\begin{figure}[H]
\centering
\begin{tikzpicture}[scale=1]
  \draw[bandac] (0,3.4) -- (3,3.4) to[out=0,in=180] (5,2.4) -- (9,2.4) node[right,font=\small] {$\Ec$};
  \draw[bandav] (0,1.4) -- (3,1.4) to[out=0,in=180] (5,0.4) -- (9,0.4) node[right,font=\small] {$\Ev$};
  \draw[fermi] (0,1.6) -- (4.6,1.6) node[above,font=\scriptsize,black,pos=0.15] {$E_{Fp}$};
  \draw[fermi] (3.4,2.2) -- (9,2.2) node[below,font=\scriptsize,black,pos=0.85] {$E_{Fn}$};
  \draw[<->,gris] (9.9,1.6) -- node[right,font=\small,black] {$qV_A$} (9.9,2.2);
  \draw[gris,dotted] (4.6,1.6) -- (9.9,1.6) (9,2.2) -- (9.9,2.2);
  % electrones inyectados desde n hacia p
  \foreach \x in {6.2,6.7,7.3}{\node[elec] at (\x,2.48) {};}
  \draw[flecha,colElec] (5.9,2.75) -- (4.2,3.2);
  \node[elec] at (2.3,3.4) {};
  \draw[flecha] (2.3,3.3) -- (2.3,1.52);
  \node[hueco] at (2.3,1.4) {};
  \draw[foton] (2.15,2.4) -- (0.6,2.4);
  % huecos inyectados desde p hacia n
  \foreach \x in {1.0,1.5,0.5}{\node[hueco] at (\x,1.3) {};}
  \draw[flecha,colHueco] (3.1,1.05) -- (4.9,0.6);
  \node[elec] at (6.6,2.4) {};
  \draw[flecha] (6.0,2.3) -- (6.0,0.52);
  \node[hueco] at (6.0,0.4) {};
  \node[elec] at (6.0,2.4) {};
  \draw[foton] (6.15,1.4) -- (7.8,1.4);
  \node[colFoton,font=\small] at (8.2,1.4) {$h\nu$};
  \node[colFoton,font=\small] at (0.3,2.4) {$h\nu$};
  \node[font=\small] at (1.2,0.2) {p}; \node[font=\small] at (8,0) {n};
\end{tikzpicture}
\caption{LED en polarización directa (Pierret, fig.~9.13). La barrera se reduce en $qV_A$;
los electrones inyectados en el lado p y los huecos inyectados en el lado n se recombinan
emitiendo fotones de energía $\approx\EG$.}
\label{fig:led-bandas}
\end{figure}

\begin{clave}[Longitud de onda emitida]
\[
  h\nu \approx \EG \quad\Longrightarrow\quad \lambda\,[\si{\micro m}]\approx\frac{1.24}{\EG\,[\si{\eV}]}
\]
El color lo fija el \textbf{gap} del material. Como electrones y huecos se distribuyen
con energías del orden de $kT$ alrededor de los bordes de banda, la emisión no es
monocromática: el pico está cerca de $\EG$ y el ancho espectral es de algunos $kT$
(típicamente \SIrange{20}{40}{nm}). Las impurezas que crean niveles dentro del gap pueden
correr el pico por debajo de $\EG$.
\end{clave}

\begin{intuicion}[por qué no hay LEDs de silicio]
En un material de gap indirecto (Si, Ge) el electrón en el mínimo de la BC y el hueco en el
máximo de la BV tienen momentos cristalinos distintos (figura~\ref{fig:Ek}). La transición
radiativa necesita un fonón y es muy lenta; antes ocurre la recombinación por centros R--G,
que entrega la energía como calor. En un material directo (GaAs, InGaN) gran parte de los
portadores se recombina radiativamente.
\end{intuicion}

\subsection{Materiales y colores}

Un LED visible necesita, según Pierret, tres cosas:
\begin{enumerate}[nosep]
  \item gap en el rango visible: $\SI{0.4}{\micro m}<\lambda<\SI{0.7}{\micro m}$, o sea
  $\SI{1.77}{\eV}<\EG<\SI{3.1}{\eV}$;
  \item gap \textbf{directo} (o un mecanismo que lo imite);
  \item que se pueda fabricar una juntura PN.
\end{enumerate}
Casi ningún semiconductor simple cumple las tres (el GaAs tiene un gap demasiado chico, el
GaP y el SiC son indirectos, a muchos II-VI no se los puede dopar tipo p). La solución
son las \textbf{aleaciones} y los \textbf{centros isoelectrónicos}:

\begin{itemize}
  \item \textbf{GaAs$_{1-x}$P$_x$}: variando $x$ el gap pasa del GaAs al GaP. Es directo
  para $x<0.45$; con $x=0.4$ emite rojo (\SI{0.65}{\micro m}). Fueron los primeros LEDs
  comerciales (displays de calculadoras de los 70).
  \item \textbf{GaP:N, GaAsP:N} (verde, amarillo, naranja): el nitrógeno reemplaza al P (es de
  la misma columna, no dona portadores: es \emph{isoelectrónico}) y crea una trampa de
  electrones $\sim\SI{0.1}{\eV}$ debajo de $\Ec$. El electrón atrapado está muy
  \emph{localizado} en el espacio, así que (por el principio de incertidumbre) su momento
  queda muy indeterminado: puede recombinarse con un hueco sin necesidad de fonón. El
  material se vuelve ``pseudo-directo''. Como el fotón sale con energía menor que $\EG$, el
  resto del cristal casi no lo reabsorbe.
  \item \textbf{AlGaAs, AlInGaP}: aleaciones de gap directo con la misma constante de red
  que el GaAs; rojos, naranjas y amarillos de alto brillo.
  \item \textbf{InGaN/GaN}: azul y verde. El GaN ($\EG=\SI{3.4}{\eV}$) es directo; la emisión
  se produce en capas de InGaN de menor gap. El LED azul eficiente (década de 1990, Premio
  Nobel 2014) completó los tres colores primarios e hizo posible el LED blanco.
\end{itemize}

\begin{table}[H]
\centering\small
\caption{LEDs comerciales (adaptado de Pierret, tabla~9.1, con agregados actuales).}
\label{tab:leds}
\begin{tabular}{llcl}
\toprule
Material & Color & $\lambda_{pico}$ (\si{\micro m}) & Comentario \\
\midrule
GaAs                 & infrarrojo  & 0.87  & controles remotos, optoacopladores \\
GaAs$_{0.6}$P$_{0.4}$  & rojo        & 0.65  & directo, baja eficiencia (\SI{0.2}{\percent}) \\
GaAs$_{0.35}$P$_{0.65}$:N & naranja  & 0.63  & centro isoelectrónico N \\
GaP:Zn--O            & rojo        & 0.70  & par Zn--O, trampa a \SI{0.3}{\eV} de $\Ec$ \\
GaP:N                & verde       & 0.565 & centro isoelectrónico N \\
AlGaAs               & rojo        & 0.65  & ``alto brillo'', heteroestructura \\
AlInGaP              & rojo a amarillo & 0.57--0.63 & alta eficiencia \\
InGaN/GaN            & azul, verde & 0.45--0.53 & los más eficientes hoy \\
InGaN + fósforo      & blanco      & --    & azul + emisión amarilla del fósforo \\
\bottomrule
\end{tabular}
\end{table}

\begin{ojo}[``LED blanco'' \diapo{23}]
No existe un semiconductor que emita blanco: el blanco es una mezcla de colores. El LED
blanco común es un LED \textbf{azul} de InGaN cubierto por un \textbf{fósforo} (típicamente
YAG:Ce) que absorbe parte del azul y reemite en el amarillo; la mezcla de azul y amarillo
se percibe como blanca. La otra opción es combinar tres LEDs (rojo, verde y azul).
\end{ojo}

\subsection{Característica eléctrica y encapsulado}

Eléctricamente el LED es un diodo: $I\approx I_0e^{V/(nV_T)}$. Lo particular es la
\textbf{tensión de encendido}, aproximadamente $\EG/q$: cada electrón que atraviesa el
diodo tiene que aportar al menos la energía del fotón que se emite ($qV\gtrsim h\nu\approx
\EG$), y además, como $I_0\propto n_i^2\propto e^{-\EG/kT}$, cuanto mayor el gap más tensión
hace falta para la misma corriente.

\begin{figure}[H]
\centering
\begin{tikzpicture}
\begin{axis}[
  width=0.65\textwidth,height=5.4cm,
  xmin=0.8,xmax=3.6,ymin=0,ymax=50,
  xlabel={$V$ (V)},ylabel={$I$ (mA)},
  axis lines=left,grid=major,grid style={gray!25},
  every axis plot/.append style={line width=1.2pt,samples=150},
]
\addplot[gris,domain=0.8:1.31] {20*exp((x-1.25)/0.06)};
\addplot[rojo,domain=1.4:2.0] {20*exp((x-1.95)/0.07)};
\addplot[verde!70!black,domain=2.0:2.47] {20*exp((x-2.4)/0.08)};
\addplot[azul,domain=2.5:3.17] {20*exp((x-3.1)/0.09)};
\node[gris,font=\small] at (axis cs:1.05,40) {IR};
\node[rojo,font=\small,anchor=east] at (axis cs:1.86,40) {rojo};
\node[verde!70!black,font=\small,anchor=east] at (axis cs:2.29,10) {amarillo};
\node[azul,font=\small,anchor=east] at (axis cs:3.0,40) {azul};
\end{axis}
\end{tikzpicture}
\caption{Curvas I--V aproximadas de LEDs de distintos colores: la tensión de encendido
crece con $\EG$ (IR GaAs $\sim\SI{1.2}{V}$, rojo $\sim\SI{1.8}{V}$, azul InGaN
$\sim\SI{3}{V}$). Compárese con el rojo de la \diapo{25}, que empieza a conducir a
\SI{1.7}{V}.}
\label{fig:led-iv}
\end{figure}

Como la corriente crece exponencialmente con la tensión, un LED \textbf{siempre} se
alimenta con limitación de corriente (una resistencia en serie o una fuente de corriente):
conectarlo directamente a una fuente de tensión lo destruye.

\paragraph{Encapsulado \diapo{23--24}.} El chip (de unos \SI{250}{\micro m} de lado) se
monta en una \emph{copa reflectora} que redirige la luz que sale por los costados; un hilo
lo conecta al otro terminal, y todo se recubre de epoxi con forma de domo. El
\textbf{ánodo} (lado p) es la pata más larga; el \textbf{cátodo} (lado n) se identifica por
el lado plano del encapsulado.

\begin{ojo}[\diapo{24}]
El dibujo indica ``N-type layer: \emph{negatively charged}''. Un semiconductor tipo n es
\textbf{eléctricamente neutro}: tiene electrones libres (negativos) pero exactamente la misma
cantidad de donores ionizados (positivos, fijos). Lo mismo vale para el tipo p (``holes for
electrons''). Sólo hay carga neta en la zona de carga espacial de la juntura.
\end{ojo}

\subsection{Heteroestructuras: confinamiento de portadores y de luz}

Los LEDs de alto brillo (y todos los láseres) usan una \textbf{doble heteroestructura}
\diapo{34}: una capa activa delgada (\SIrange{0.1}{0.5}{\micro m}) de gap chico, por
ejemplo GaAs o AlGaAs con poco Al, entre dos capas de \emph{confinamiento} de gap mayor
(Al$_x$Ga$_{1-x}$As con más Al), una n y otra p.

\begin{figure}[H]
\centering
\begin{tikzpicture}[scale=1]
  % bandas en directa (simplificado, sin curvaturas)
  \draw[bandac] (0,3.0) -- (3,3.0) -- (3,2.5) -- (5,2.5) -- (5,3.6) -- (8,3.6);
  \draw[bandav] (0,0.6) -- (3,0.6) -- (3,1.1) -- (5,1.1) -- (5,0.0) -- (8,0.0);
  \node[colEc,font=\small,anchor=west] at (8,3.6) {$\Ec$};
  \node[colEv,font=\small,anchor=west] at (8,0.0) {$\Ev$};
  \fill[pattern=north east lines,pattern color=azul!40] (3,2.5) rectangle (5,3.0);
  \fill[pattern=north east lines,pattern color=rojo!40] (3,0.6) rectangle (5,1.1);
  \foreach \x in {3.3,3.8,4.3,4.7}{\node[elec] at (\x,2.58) {};}
  \foreach \x in {3.4,3.9,4.4}{\node[hueco] at (\x,1.02) {};}
  \draw[flecha,colElec] (1,3.12) -- (2.8,2.75);
  \node[colElec,font=\scriptsize,align=center] at (1.2,3.55) {electrones\\inyectados};
  \draw[flecha,colHueco] (7,0.12) -- (5.2,0.9);
  \node[colHueco,font=\scriptsize,align=center] at (6.9,0.65) {huecos\\inyectados};
  \node[font=\scriptsize,align=center] at (6.4,3.05) {barrera para\\electrones};
  \node[font=\scriptsize,align=center] at (1.6,1.15) {barrera para\\huecos};
  \draw[foton] (3.9,1.8) -- (2.0,1.8);
  \node[colFoton,font=\small] at (1.75,1.8) {$h\nu$};
  \draw[dashed,gris] (3,-0.4) -- (3,4.0) (5,-0.4) -- (5,4.0);
  \node[font=\scriptsize,align=center] at (1.5,-0.7) {n-AlGaAs\\(gap grande)};
  \node[font=\scriptsize,align=center] at (4,-0.7) {capa activa\\(gap chico)};
  \node[font=\scriptsize,align=center] at (6.5,-0.7) {p-AlGaAs\\(gap grande)};
  % indice de refraccion
  \begin{scope}[yshift=-3.1cm]
    \draw[flecha] (0,0) -- (8.3,0) node[right,font=\small] {$x$};
    \draw[flecha] (0,0) -- (0,1.3) node[above,font=\small] {$n_{\text{refr}}$};
    \draw[azul,line width=1.1pt] (0,0.5) -- (3,0.5) -- (3,1.0) -- (5,1.0) -- (5,0.5) -- (8,0.5);
    \node[font=\scriptsize,anchor=west] at (5.2,1.0) {mayor índice en la capa activa: guía de onda};
  \end{scope}
\end{tikzpicture}
\caption{Doble heteroestructura en polarización directa (simplificada, \diapo{34}). Los
escalones de banda atrapan electrones y huecos en la capa activa; el mayor índice de
refracción de esa capa guía la luz.}
\label{fig:doble-hetero}
\end{figure}

\begin{itemize}
  \item \textbf{Confinamiento de portadores:} los escalones $\Delta E_c$ y $\Delta E_v$ son
  barreras que impiden que electrones y huecos escapen de la capa activa. Al quedar
  concentrados en un volumen chico, su densidad es alta y la recombinación radiativa (que
  crece con el producto $np$) le gana a la no radiativa.
  \item \textbf{Ventanas transparentes:} las capas de confinamiento tienen gap mayor que la
  energía de los fotones emitidos, así que no los reabsorben.
  \item \textbf{Confinamiento óptico:} la capa activa tiene mayor índice de refracción y
  actúa como guía de onda (imprescindible en el láser).
\end{itemize}

\subsection{Eficiencia y extracción de la luz}

\begin{definicion}[eficiencias de un LED]
\[
  \eta_{ext} = \frac{\text{fotones que salen al exterior por segundo}}{\text{electrones que circulan por segundo}}
  = \eta_{iny}\cdot\eta_{int}\cdot\eta_{extr}
\]
\begin{itemize}[nosep]
  \item $\eta_{iny}$: fracción de la corriente que inyecta portadores en la región activa.
  \item $\eta_{int}$ (eficiencia cuántica interna): fracción de las recombinaciones que son
  radiativas, $\eta_{int}=R_{rad}/(R_{rad}+R_{no\,rad})$.
  \item $\eta_{extr}$: fracción de los fotones generados que logra salir del chip.
\end{itemize}
La \emph{eficiencia de enchufe} (\emph{wall-plug}) es $P_{luz}/(IV)$.
\end{definicion}

\paragraph{El problema de la extracción.} Los semiconductores tienen índices de refracción
altos ($n\approx3.4$ para GaP o GaAs). Por la ley de Snell, la luz que llega a la
superficie con un ángulo mayor que el \textbf{ángulo crítico}
\[
  \sin\theta_c = \frac{n_2}{n_1}
\]
sufre \emph{reflexión total interna} y queda atrapada. Promediando sobre todas las
direcciones (Pierret, ecs.~9.8--9.10), la fracción que sale es
\[
  \overline{T} \approx \frac{\sin^2\theta_c}{2}\,T_N
  = \frac{2}{n_r(1+n_r)^2},
  \qquad n_r=\frac{n_1}{n_2},
\]
donde $T_N=4n_r/(1+n_r)^2$ es la transmisión a incidencia normal.

\begin{figure}[H]
\centering
\begin{tikzpicture}[scale=1.1]
  \fill[azulclaro] (-3,-2) rectangle (3,0);
  \draw[line width=0.8pt] (-3,0) -- (3,0);
  \node[font=\small,anchor=west] at (-2.9,-1.75) {semiconductor, $n_1=3.4$};
  \node[font=\small,anchor=west] at (-3,0.6) {aire, $n_2=1$};
  \coordinate (O) at (0,-1.6);
  \fill[naranja!35] (O) -- ({-1.6*tan(17.1)},0) -- ({1.6*tan(17.1)},0) -- cycle;
  \draw[naranja,line width=0.8pt] (O) -- ({-1.6*tan(17.1)},0) (O) -- ({1.6*tan(17.1)},0);
  \draw[dashed,gris] (O) -- (0,0.9);
  \draw[flecha,naranja!80!black] (O) -- (0.2,0) -- (0.32,0.8);
  \draw[flecha,azul] (O) -- (1.8,0) -- (3.0,-1.1);
  \node[azul,font=\scriptsize,align=center] at (2.2,-1.5) {reflexión\\total interna};
  \draw (0,-1.0) arc[start angle=90,end angle={90-17.1},radius=0.6];
  \node[font=\scriptsize] at (0.33,-0.65) {$\theta_c$};
  \node[elec,fill=colFoton] at (O) {};
  \node[font=\scriptsize,anchor=east] at (-0.1,-1.6) {emisión};
  \node[naranja!80!black,font=\scriptsize,align=center] at (-1.2,-1.1) {cono de\\escape};
\end{tikzpicture}
\caption{Cono de escape: sólo los rayos dentro del ángulo crítico pueden salir del chip.
Para GaP--aire, $\theta_c=\ang{17.1}$.}
\label{fig:cono}
\end{figure}

\begin{ejemplo}[extracción de luz en GaP]
GaP ($n_1=3.4$) hacia aire ($n_2=1$): $T_N=4\cdot3.4/4.4^2=\SI{70}{\percent}$ a incidencia
normal, pero $\theta_c=\arcsin(1/3.4)=\ang{17.1}$ y $\overline{T}=2/(3.4\cdot4.4^2)
\approx\SI{3}{\percent}$. ¡Sólo sale el \SI{3}{\percent} de la luz generada!

Con epoxi ($n_2=1.5$): $n_r=2.27$, $\theta_c=\ang{26.2}$ y
$\overline{T}=2/(2.27\cdot3.27^2)\approx\SI{8}{\percent}$: casi el triple. El domo del
encapsulado hace que luego la luz llegue a la interfaz epoxi--aire casi en forma normal y
salga. Por eso los LEDs se encapsulan en epoxi en forma de domo, se usan sustratos
transparentes y espejos traseros, y se texturizan o se dan formas especiales a los chips.
\end{ejemplo}

\subsection{Complemento: el diodo láser}

Si la inyección es tan grande que en la capa activa hay más electrones en la BC que en la
BV (\emph{inversión de población}), un fotón que pasa estimula la emisión de otro
\emph{idéntico} (misma $\lambda$, fase y dirección): \textbf{emisión estimulada}. Con caras
del cristal que actúan como espejos (cavidad) y por encima de una \emph{corriente umbral},
la ganancia supera a las pérdidas y el dispositivo emite luz coherente, muy
monocromática y direccional. La doble heteroestructura es la que hace posible alcanzar la
inversión con corrientes razonables a temperatura ambiente.

% =====================================================================
\section{Resumen y repaso}
% =====================================================================

\subsection{Comparación de dispositivos}

\begin{table}[H]
\centering\small
\caption{Los dispositivos de la unidad de un vistazo.}
\label{tab:comparacion}
\begin{tabularx}{\textwidth}{l>{\raggedright\arraybackslash}p{2.6cm}c>{\raggedright\arraybackslash}X>{\raggedright\arraybackslash}X}
\toprule
Dispositivo & Polarización & Cuadrante & Figura de mérito & Limitación principal \\
\midrule
Fotoconductor & cualquiera (sin juntura) & -- & ganancia $G_{opt}=\tau/t_{tr}$ & lento ($\sim\tau$) \\
Fotodiodo PN  & inversa & III & $\eta$, $\Resp$ & difusión: decenas de MHz \\
Fotodiodo PIN & inversa & III & $\eta$, $\Resp$, $f_{\max}=\vsat/W_i$ & compromiso $\eta$--velocidad \\
APD           & inversa, cerca de ruptura & III & ganancia $M$ & ruido de multiplicación, alta tensión \\
Celda solar   & ninguna (carga) & IV & $\eta=P_m/P_{in}$, $FF$, $\Voc$, $\Isc$ & termalización, transparencia, $R_s$ \\
LED           & directa & I & $\eta_{ext}$, $\lambda\approx1.24/\EG$ & extracción, gap directo \\
\bottomrule
\end{tabularx}
\end{table}

\subsection{Fórmulas clave}

\begin{tcolorbox}[base,colback=white,colframe=verde,title={Formulario}]
\renewcommand{\arraystretch}{1.6}
\begin{tabularx}{\linewidth}{>{\raggedright\arraybackslash}p{0.33\linewidth}X}
Energía del fotón & $\Eph=h\nu=\dfrac{hc}{\lambda}$,\quad $\lambda[\si{\micro m}]=\dfrac{1.24}{\Eph[\si{\eV}]}$,\quad $\lG=\dfrac{1.24}{\EG}$ \\
Absorción (Beer--Lambert) & $\Phi(x)=\Phi_0e^{-\alpha x}$,\quad penetración $1/\alpha$ \\
Fotogeneración & $G_L(x)=\alpha\Phi_0e^{-\alpha x}$,\quad $\Phi_0=\dfrac{(1-R)P}{h\nu A}$ \\
Reflexión normal & $R=\left(\dfrac{n-1}{n+1}\right)^2$ \\
Exceso estacionario & $\Delta n=\Delta p=G_L\tau$ \\
Fotoconductor & $\Delta\sigma=q(\mu_n+\mu_p)\Delta n$,\quad $G_{opt}=\dfrac{\tau_n}{t_{tr}}\left(1+\dfrac{\mu_p}{\mu_n}\right)$,\quad $t_{tr}=\dfrac{L^2}{\mu_nV}$ \\
Fotodiodo PN & $I=I_0(e^{V_A/V_T}-1)-I_L$,\quad $I_L=qAG_L(L_n+W+L_p)$ \\
Eficiencia cuántica & $\eta=\dfrac{I_L/q}{\Popt/h\nu}$;\quad PIN: $\eta=\eta_i(1-R)(1-e^{-\alpha W})$ \\
Responsividad & $\Resp=\dfrac{I_L}{\Popt}=\dfrac{\eta q}{h\nu}=\eta\dfrac{\lambda[\si{\micro m}]}{1.24}$ [\si{A/W}] \\
Velocidad del PIN & $f_{\max}\approx\vsat/W_i$,\quad $C_j=\epsilon A/W_i$ \\
Celda solar & $\Isc=I_L$,\quad $\Voc=V_T\ln(1+I_L/I_0)$ \\
 & $FF=\dfrac{V_mI_m}{\Isc\Voc}$,\quad $\eta=\dfrac{FF\,\Isc\Voc}{P_{in}}$ \\
LED & $\lambda\approx1.24/\EG$,\quad $\eta_{ext}=\eta_{iny}\eta_{int}\eta_{extr}$ \\
Extracción & $\sin\theta_c=n_2/n_1$,\quad $\overline{T}\approx\dfrac{2}{n_r(1+n_r)^2}$ \\
\end{tabularx}
\end{tcolorbox}

\subsection{Preguntas de repaso}

\begin{enumerate}
  \item \textbf{¿Por qué un fotodiodo de Si no detecta \SI{1.55}{\micro m}?}\\
  Porque $\Eph=\SI{0.8}{\eV}<\EG=\SI{1.12}{\eV}$: el Si es transparente a esa luz.
  \item \textbf{¿Por qué la responsividad ideal crece linealmente con $\lambda$?}\\
  Porque a igual potencia llegan más fotones cuanto menor es su energía, y cada fotón
  aporta (como mucho) un electrón: $\Resp=\eta\lambda/1.24$.
  \item \textbf{¿Por qué el fotodiodo PN es lento y el PIN rápido?}\\
  En el PN la mayoría de los pares se genera en zonas neutras y llega a la juntura por
  difusión (ns). En el PIN se genera en la capa i, donde hay campo, y se colecta por deriva
  a $\vsat$ (ps).
  \item \textbf{¿Qué pasa con $\Voc$ e $\Isc$ de una celda si se duplica la luz?}\\
  $\Isc$ se duplica; $\Voc$ aumenta sólo $V_T\ln2\approx\SI{18}{mV}$.
  \item \textbf{¿Y si sube la temperatura?}\\
  $I_0$ crece mucho (como $n_i^2$), $\Voc$ baja (unos \SI{2}{mV/\celsius} en Si) y con ella la
  potencia, aunque $\Isc$ aumente levemente (baja $\EG$).
  \item \textbf{¿Por qué un gap muy chico o muy grande es malo para una celda?}\\
  Muy chico: se absorbe casi todo pero cada fotón aporta poco ($\EG$) y el resto se pierde
  como calor; además $\Voc$ es chica. Muy grande: se aprovecha bien cada fotón absorbido,
  pero la mayoría de los fotones no se absorbe.
  \item \textbf{¿Cómo puede un fotoconductor tener ganancia mayor que 1 y un fotodiodo no?}\\
  En el fotoconductor los contactos óhmicos reponen portadores: un electrón circula
  $\tau/t_{tr}$ veces antes de que su par se recombine. En el fotodiodo cada par se
  separa una sola vez (salvo multiplicación por avalancha).
  \item \textbf{¿Por qué los LEDs se hacen con materiales de gap directo, y cómo emite
  luz el GaP, que es indirecto?}\\
  En gap directo la recombinación radiativa es rápida y le gana a la no radiativa. En GaP
  se agregan centros isoelectrónicos (N, Zn--O) que localizan al electrón y relajan la
  conservación del momento.
  \item \textbf{¿Para qué sirve la doble heteroestructura?}\\
  Confina portadores en la capa activa (más recombinación radiativa), evita la
  reabsorción (capas de gap mayor) y guía la luz (mayor índice).
\end{enumerate}

\subsection{Correcciones a las diapositivas (resumen)}

\begin{table}[H]
\centering\small
\begin{tabularx}{\textwidth}{cX}
\toprule
Diapo & Corrección \\
\midrule
6  & GaP es de gap \emph{indirecto} (no emite verde por sí solo). \\
9  & Al bajar $\lambda$ no aumenta $\EG$: aumenta la energía del fotón (menos fotones por watt). \\
10 & $\vsat\approx\SI{e7}{cm/s}$, no \SI{e5}{cm/s}. El PIN es una estructura, no un modo. ``Longitud de deplexión'' $\to$ ancho de la ZCE. \\
11--12 & El lado p de una homojuntura no es transparente: se hace delgado. La absorción no ocurre sólo en la ZCE. \\
19 & La corriente de oscuridad se filtra en continua, pero su ruido de granalla no. \\
20 & $\Resp$ es por unidad de \emph{potencia} [\si{A/W}]; la ``frecuencia'' es la de la luz. \\
24 & Las capas n y p son eléctricamente neutras. \\
27 & 30--35\,\% corresponde a celdas III-V multijuntura; el gap del Si está cerca del óptimo. \\
38 & La celda partida reduce pérdidas óhmicas ($I^2R$), no ``térmicas''. \\
44 & $\Delta R\neq L/(\Delta\sigma A)$: lo aditivo es la conductancia. \\
46 & Exacto: $G_{opt}=\tau_n(\mu_n+\mu_p)\mathcal{E}/L$; la diapo desprecia $\mu_p$. \\
\bottomrule
\end{tabularx}
\end{table}

\subsection*{Bibliografía}
\begin{itemize}[nosep]
  \item R.~F.~Pierret, \emph{Semiconductor Device Fundamentals}, Addison-Wesley, 1996:
  §3.3 (R--G, fotogeneración) y cap.~9 (diodos optoelectrónicos).
  \item A.~S.~Sedra, K.~C.~Smith, \emph{Microelectronic Circuits}, 7.\textsuperscript{a} ed.,
  Oxford, 2015: cap.~15.
  \item R.~N.~Selva, \emph{Dispositivos Electrónicos}, 2.\textsuperscript{a} ed., Nueva
  Librería, 2008: cap.~15.1 (fotoconductores).
  \item PVEducation: \texttt{https://www.pveducation.org} (celdas solares).
  \item Diapositivas de clase ``Física Electrónica -- Optoelectrónica'', Ing.~E.~D.~Baez, ITBA, 2026.
\end{itemize}


\end{document}
